\documentclass[12pt,a4paper]{report}

% ============================================================
% PACKAGES
% ============================================================
\usepackage[a4paper, top=1in, bottom=1in, left=1.5in, right=1in]{geometry}
\usepackage{setspace}
\usepackage[T1]{fontenc}
\usepackage[utf8]{inputenc}
\usepackage{mathptmx}
\usepackage{amsmath}
\usepackage{amssymb}
\usepackage{booktabs}
\usepackage{tabularx}
\usepackage{array}
\usepackage{graphicx}
\usepackage{caption}
\usepackage{subcaption}
\usepackage{fancyhdr}
\usepackage{titlesec}
\usepackage{tocloft}
\usepackage[hidelinks,hypertexnames=false]{hyperref}
\usepackage{url}
\usepackage{enumitem}
\usepackage{parskip}
\usepackage{microtype}
\usepackage{xcolor}
\usepackage{listings}
\usepackage{float}
\usepackage{ragged2e}
\usepackage{tikz}
\usepackage{pgfplots}
\pgfplotsset{compat=1.18}
\usetikzlibrary{shapes.geometric, arrows.meta, positioning, fit,
                backgrounds, calc, decorations.pathreplacing, shadows,
                shapes.multipart, matrix}

% ============================================================
% COLOUR PALETTE (AWS-inspired for architecture diagrams)
% ============================================================
\definecolor{awsorange}{RGB}{255,153,0}
\definecolor{awsblue}{RGB}{35,103,178}
\definecolor{awsgreen}{RGB}{59,153,75}
\definecolor{awspurple}{RGB}{142,68,173}
\definecolor{awsgray}{RGB}{100,100,100}
\definecolor{awslightblue}{RGB}{213,232,255}
\definecolor{awslightgreen}{RGB}{213,245,218}
\definecolor{awslightorange}{RGB}{255,235,200}
\definecolor{awslightpurple}{RGB}{235,218,245}
\definecolor{sectionbg}{RGB}{245,245,245}
\definecolor{propblue}{RGB}{41,128,185}
\definecolor{propgreen}{RGB}{39,174,96}
\definecolor{propred}{RGB}{192,57,43}
\definecolor{propgray}{RGB}{149,165,166}

% ============================================================
% AIT FORMATTING
% ============================================================
\onehalfspacing
\setlength{\parindent}{0pt}
\setlength{\parskip}{8pt}

% ---- Page headers via fancyhdr ----
\setlength{\headheight}{14.5pt}
\pagestyle{fancy}
\fancyhf{}
\fancyhead[L]{\small\leftmark}
\fancyhead[R]{\small\thepage}
\fancyfoot{}
\renewcommand{\headrulewidth}{0.4pt}
\fancypagestyle{plain}{%
  \fancyhf{}
  \fancyfoot[C]{\small\thepage}
  \renewcommand{\headrulewidth}{0pt}
}

\titleformat{\chapter}[display]
  {\normalfont\fontsize{14}{16}\bfseries\centering}
  {\MakeUppercase{\chaptertitlename}\ \thechapter}{0pt}{\MakeUppercase}
\titlespacing*{\chapter}{0pt}{0pt}{20pt}

\titleformat{name=\chapter,numberless}[block]
  {\normalfont\fontsize{14}{16}\bfseries\centering}
  {}{0pt}{\MakeUppercase}
\titlespacing*{name=\chapter,numberless}{0pt}{0pt}{20pt}

\titleformat{\section}
  {\normalfont\fontsize{12}{14}\bfseries}{\thesection}{1em}{}
\titlespacing*{\section}{0pt}{14pt}{6pt}

\titleformat{\subsection}
  {\normalfont\fontsize{12}{14}\bfseries}{\thesubsection}{1em}{}
\titlespacing*{\subsection}{0pt}{12pt}{4pt}

\titleformat{\subsubsection}
  {\normalfont\fontsize{12}{14}\bfseries}{\thesubsubsection}{1em}{}
\titlespacing*{\subsubsection}{0pt}{10pt}{4pt}

\renewcommand{\cfttoctitlefont}{\normalfont\fontsize{14}{16}\bfseries}
\renewcommand{\contentsname}{CONTENTS}
\renewcommand{\listfigurename}{LIST OF FIGURES}
\renewcommand{\listtablename}{LIST OF TABLES}
\setlength{\cftbeforetoctitleskip}{0pt}
\setlength{\cftaftertoctitleskip}{14pt}
\renewcommand{\cftchapfont}{\bfseries}
\renewcommand{\cftchappagefont}{\bfseries}
\setlength{\cftbeforechapskip}{4pt}

\captionsetup[table]{position=above, labelsep=period,
  justification=raggedright, singlelinecheck=false, font=normalsize}
\captionsetup[figure]{position=below, labelsep=period,
  justification=justified, font=normalsize, singlelinecheck=false}

% Prevent overfull hbox in general text
\sloppy

\lstset{
  basicstyle=\ttfamily\small, breaklines=true, frame=single,
  backgroundcolor=\color{gray!8},
  xleftmargin=10pt, xrightmargin=10pt, aboveskip=8pt, belowskip=8pt
}

\graphicspath{{figures/}}

% ============================================================
% TIKZ STYLES
% ============================================================
\tikzset{
  % Pipeline boxes
  pbox/.style={rectangle, rounded corners=4pt, draw=#1!70, fill=#1!12,
               text centered, font=\small\bfseries, minimum height=0.65cm,
               minimum width=2.4cm, text width=2.2cm, drop shadow},
  % Arrow
  parrow/.style={-{Stealth[length=6pt]}, thick, color=awsgray},
  % Layer label
  layerlabel/.style={font=\footnotesize\bfseries, color=#1},
  % Component node
  compnode/.style={rectangle, rounded corners=3pt, draw=#1, fill=white,
                   font=\scriptsize, minimum height=0.55cm,
                   minimum width=1.8cm, text width=1.7cm, text centered},
  % Dashed bounding box
  boundbox/.style={rectangle, rounded corners=6pt, draw=#1!60, fill=#1!5,
                   dashed, thick},
  % Decision diamond
  decision/.style={diamond, draw=awsblue!70, fill=awslightblue,
                   font=\scriptsize, aspect=2.2, text centered,
                   minimum height=0.7cm, minimum width=2.0cm},
  % Evaluation row
  evalbox/.style={rectangle, rounded corners=3pt, draw=#1!70, fill=#1!10,
                  font=\small, minimum height=0.6cm, text width=3.5cm,
                  text centered},
}

% ============================================================
% DOCUMENT
% ============================================================
\begin{document}
\pagenumbering{roman}

% ============================================================
% TITLE PAGE
% ============================================================
\begin{titlepage}
\begin{center}
\vspace*{0.5cm}
{\fontsize{14}{18}\bfseries\MakeUppercase{%
TechPulse: A Real-Time Hybrid Retrieval-Augmented Generation\\
System for Emerging Technology Intelligence}\\[4pt]
\normalsize\mdseries\textit{Progress Report}}

\vspace{1.8cm}
{\normalsize by}

    \vspace{1em}

    \begin{tabular}{cc}
        Aye Khin Khin Hpone (Yolanda Lim) & Dechathon Niamsa-Ard \\
        {\small st125970} & {\small st126235}
    \end{tabular}

\vfill

{\normalsize AT82.9002 Selected Topic: Data Engineering and MLOps}\\[6pt]
{\normalsize Instructor: Dr.\ Chantri Polprasert}\\[4pt]
{\normalsize Teaching Assistants: Swaraj Bhanja, Suryansh Srivastava}

\vfill
{\normalsize
Asian Institute of Technology\\[4pt]
School of Engineering and Technology\\[4pt]
Thailand\\[8pt]
March 2026
}
\end{center}
\end{titlepage}

% ============================================================
% ABSTRACT
% ============================================================
\chapter*{Abstract}
\addcontentsline{toc}{chapter}{ABSTRACT}

The rapid expansion of emerging technologies has created a fragmented knowledge ecosystem where scholarly research coexists with fast-moving industry discussions. Large language models (LLMs) offer advanced generative capabilities but are constrained by static training corpora, producing hallucinated or outdated outputs when queried about recent developments. Retrieval-Augmented Generation (RAG) mitigates this by conditioning generation on externally retrieved evidence; however, baseline vector-only retrieval often suffers from contextual imprecision.

This report presents the implementation progress of TechPulse, a cloud-native Hybrid RAG system designed to synthesize emerging technology intelligence under operational latency and cost constraints. The system ingests data from five heterogeneous sources---ArXiv, Hacker News, DEV.to, GitHub Trending, and curated RSS/Atom feeds---with deterministic preprocessing, token-based chunking, and semantic embedding via the MiniLM sentence-transformer. Embeddings are indexed in PostgreSQL with the pgvector extension, supporting integrated relational filtering and vector similarity search.

The retrieval pipeline operates as a three-stage process: metadata-constrained filtering, cosine similarity search, and lightweight reranking incorporating semantic similarity, lexical overlap, and exponential recency weighting. An automatic LLM backend fallback chain ensures generation availability across deployment environments, with Groq (\texttt{llama-3.1-8b-instant}) as the primary backend under AWS Academy Learner Lab constraints. The system further incorporates dual-criteria drift detection (Shewhart $3\sigma$ control charts and relative-drop thresholds), three-layer hallucination verification, and comprehensive health monitoring. A controlled evaluation framework comparing baseline and hybrid retrieval has been developed using RAGAS metrics and grid search over reranking weights.

This report documents the system architecture, pipeline design, implementation decisions, deviations from the original proposal, current progress, and remaining work.

\medskip
\noindent\textbf{Keywords:} Retrieval-Augmented Generation; Hybrid Information Retrieval; Vector Databases; Cloud-Native Architecture; MLOps; Semantic Embedding; Emerging Technology Intelligence.

% ============================================================
% ACKNOWLEDGMENTS
% ============================================================
\chapter*{Acknowledgments}
\addcontentsline{toc}{chapter}{ACKNOWLEDGMENTS}

The authors express sincere gratitude to Dr.\ Chantri Polprasert, the course instructor, for providing guidance on the system design and evaluation methodology throughout this project. We also thank the teaching assistants, Mr.\ Swaraj Bhanja and Mr.\ Suryansh Srivastava, for their constructive feedback during development milestones.

This project was developed using the AWS Academy Learner Lab environment provided by the Asian Institute of Technology (AIT). We acknowledge AWS Academy for providing cloud computing resources that enabled the deployment and evaluation of the TechPulse system.

% ============================================================
% FRONT MATTER
% ============================================================
\newpage\tableofcontents
\newpage\addcontentsline{toc}{chapter}{LIST OF TABLES}\listoftables
\newpage\addcontentsline{toc}{chapter}{LIST OF FIGURES}\listoffigures

% ============================================================
% LIST OF ABBREVIATIONS
% ============================================================
\newpage
\chapter*{List of Abbreviations}
\addcontentsline{toc}{chapter}{LIST OF ABBREVIATIONS}

\begin{tabular}{@{}p{3.2cm}p{10.0cm}@{}}
\textbf{AIT} & Asian Institute of Technology \\
\textbf{ANN} & Approximate Nearest Neighbor \\
\textbf{API} & Application Programming Interface \\
\textbf{CI/CD} & Continuous Integration / Continuous Deployment \\
\textbf{CORS} & Cross-Origin Resource Sharing \\
\textbf{DLQ} & Dead Letter Queue \\
\textbf{HF} & HuggingFace \\
\textbf{HNSW} & Hierarchical Navigable Small World \\
\textbf{IaC} & Infrastructure as Code \\
\textbf{LCL} & Lower Control Limit \\
\textbf{LLM} & Large Language Model \\
\textbf{MLOps} & Machine Learning Operations \\
\textbf{NLP} & Natural Language Processing \\
\textbf{RAG} & Retrieval-Augmented Generation \\
\textbf{RAGAS} & Retrieval-Augmented Generation Assessment \\
\textbf{RDS} & Relational Database Service \\
\textbf{RSS} & Really Simple Syndication \\
\textbf{S3} & Simple Storage Service \\
\textbf{SAM} & Serverless Application Model \\
\textbf{SHA-256} & Secure Hash Algorithm (256-bit) \\
\textbf{SNS} & Simple Notification Service \\
\textbf{SPA} & Single-Page Application \\
\textbf{SQL} & Structured Query Language \\
\textbf{SQS} & Simple Queue Service \\
\textbf{VPC} & Virtual Private Cloud \\
\end{tabular}

\newpage
\pagenumbering{arabic}

% ============================================================
% CHAPTER 1: INTRODUCTION
% ============================================================
\chapter{Introduction}

\section{Problem Description and Motivation}
\label{sec:problem}

\textbf{Motivating scenario.} A machine learning engineer at a technology company needs to stay current with advances in transformer efficiency. She must manually monitor ArXiv preprints, scan Hacker News threads, and cross-reference multiple sources---a process that is time-consuming, fragmented, and prone to missing recent developments. An intelligent assistant capable of synthesizing both scholarly and community-sourced knowledge into a single, citation-grounded response would significantly accelerate her workflow. TechPulse is designed precisely for this use case.

The exponential growth of emerging technologies---including artificial intelligence, machine learning, distributed systems, and advanced computing---has produced an unprecedented volume of technical knowledge distributed across heterogeneous platforms. Scholarly repositories such as ArXiv provide structured, peer-reviewed or preprint research outputs contributing to theoretical and methodological advancement. Community-driven platforms and industry discussion forums such as Hacker News reflect near-real-time practitioner experiences and emerging tool ecosystems. However, the fragmentation of these knowledge regimes creates an information landscape that is increasingly difficult for technology professionals to synthesize effectively and in a timely manner.

Purely generative large language model (LLM) assistants are insufficient for emerging technology intelligence tasks due to their static knowledge cutoff, susceptibility to hallucinated responses, lack of verifiable evidence grounding, and inability to integrate with continuously updating data streams (Ji et al., 2023). Although Retrieval-Augmented Generation (RAG) architectures partially address these limitations by incorporating external context, baseline vector-only retrieval approaches may still retrieve semantically adjacent but contextually irrelevant documents. Such systems often fail to incorporate temporal relevance signals, structured metadata constraints, or recency-aware ranking mechanisms, which are essential for accurately capturing rapidly evolving technology trends.

Moreover, many existing RAG research prototypes focus primarily on retrieval accuracy without addressing production-readiness: deterministic pipeline design, idempotent ingestion, monitoring and observability, deployment automation, and cloud cost governance are frequently neglected. This limits their applicability in real-world operational environments.

The central engineering and research challenge addressed in this project is therefore:

\begin{quote}
\textbf{\textit{How can a scalable, continuously updating hybrid retrieval system be designed and deployed to synthesize emerging technology trends from heterogeneous scholarly and industry sources while satisfying latency and cost constraints?}}
\end{quote}

This project is motivated by three converging factors. First, technology professionals increasingly require timely and citation-grounded intelligence that synthesizes insights from both scholarly research and rapidly evolving industry discussions. Second, mature serverless cloud infrastructures now enable cost-effective deployment of production-grade data and machine learning pipelines. Third, a gap remains in the Retrieval-Augmented Generation (RAG) literature between retrieval-focused research prototypes and operationally robust systems capable of continuous ingestion, evaluation, and deployment in real-world environments.

\section{Research Questions}
\label{sec:rq}

To address the central problem, this study investigates the following research questions:

\begin{enumerate}[leftmargin=0.5in]
  \item \textbf{Data Engineering.} How can a cloud-native pipeline be designed to reliably ingest and preprocess heterogeneous scholarly and near-real-time industry data in a deterministic and idempotent manner?

  \item \textbf{Retrieval Optimization.} Does hybrid retrieval---combining metadata filtering, vector similarity search, and lightweight reranking---improve retrieval precision and grounding faithfulness compared with baseline vector-only retrieval?

  \item \textbf{Operational Trade-offs.} What trade-offs arise between retrieval precision, p95 latency, and cost per query under realistic workload conditions?

  \item \textbf{MLOps Integration.} How can deployment, monitoring, version control, and cost governance be integrated to ensure the reliability and maintainability of a production-grade RAG system?
\end{enumerate}

\section{Objectives of the Study}
\label{sec:objectives}

The primary objective is to design, implement, and evaluate a hybrid Retrieval-Augmented Generation system capable of synthesizing real-time emerging technology intelligence within a scalable, cost-aware cloud-native architecture. The specific objectives are:

\begin{enumerate}[leftmargin=0.5in]
  \item To construct a cloud-native data engineering pipeline for batch and near-real-time ingestion from five heterogeneous data sources.
  \item To implement and compare baseline vector-only retrieval against hybrid retrieval with metadata filtering and reranking-lite optimization.
  \item To evaluate system performance across context precision, faithfulness, p95 latency, and cost per query.
  \item To deploy and monitor the system using serverless AWS infrastructure consistent with MLOps best practices (Sculley et al., 2015; Breck et al., 2017).
\end{enumerate}

\section{Significance of the Study}
\label{sec:significance}

This project contributes at the intersection of data engineering, natural language processing, and MLOps. From a technical perspective, it demonstrates how retrieval optimization can be integrated within a production-grade cloud architecture without relying on computationally expensive cross-encoder rerankers or specialized vector database infrastructure. The system provides a practical framework for grounded intelligence synthesis across academic and industry domains while maintaining architectural simplicity suitable for scalable deployment. From a practical perspective, the resulting system offers a blueprint for citation-supported technology intelligence synthesis that is operationally maintainable, cost-governed, and deployable within resource-constrained environments such as the AWS Academy Learner Lab.

\section{Report Structure}

Chapter~1 introduces the problem description, motivation, research questions, and objectives. Chapter~2 reviews relevant literature on Retrieval-Augmented Generation, hybrid retrieval strategies, vector databases, and MLOps practices. Chapter~3 presents the implemented methodology and technical approach. Chapter~4 describes the system architecture and data engineering pipeline design as implemented. Chapter~5 outlines the evaluation framework and metrics. Chapter~6 documents implementation progress, deviations from the original proposal, and current system status. Chapter~7 summarizes remaining work, known limitations, and the plan toward final evaluation and deployment. References and appendices follow.

% ============================================================
% CHAPTER 2: BACKGROUND AND RELATED WORK
% ============================================================
\chapter{Background and Related Work}

\section{Retrieval-Augmented Generation}
\label{sec:lit_rag}

Retrieval-Augmented Generation (RAG) was introduced by Lewis et al. (2020) as a framework for augmenting sequence-to-sequence generation with a dense retrieval component, enabling the model to access external documents at inference time. This approach significantly improves performance on knowledge-intensive tasks by conditioning generation on retrieved evidence rather than relying solely on parametric memory. Despite their scale, LLMs remain prone to hallucination when responding to domain-specific or time-sensitive queries (Ji et al., 2023), making retrieval-based augmentation a practically important technique.

Recent surveys (Gao et al., 2023; Zhang et al., 2024) categorize RAG systems into four architectural types: (1) single-stage dense retrieval; (2) hybrid lexical--semantic retrieval; (3) multi-stage retrieval with reranking; and (4) iterative or self-reflective retrieval. While baseline dense retrieval improves factual grounding, empirical benchmarks such as BEIR (Thakur et al., 2021) demonstrate that semantic similarity alone is insufficient for high-precision retrieval in dynamic or heterogeneous corpora. Foundational retrieval-augmented systems such as Atlas (Izacard et al., 2022) and RETRO (Borgeaud et al., 2022) demonstrate retrieval augmentation at scale but require computationally intensive pretraining pipelines unsuited to cost-constrained production environments.

\section{Embedding Models and Vector Search}
\label{sec:lit_embed}

Dense retrieval relies on embedding models that transform text into high-dimensional vector representations. Sentence-transformer architectures such as MiniLM (Wang et al., 2020) provide computationally efficient embeddings with favorable trade-offs between semantic expressiveness and inference latency. PostgreSQL with the pgvector extension enables cosine similarity search alongside SQL-based metadata filtering within a single relational storage layer, reducing architectural complexity compared to dedicated vector databases such as FAISS or Milvus. This integrated approach is increasingly favored for production RAG deployments (Gao et al., 2023).

\section{Hybrid Retrieval and Reranking}
\label{sec:lit_hybrid}

Hybrid retrieval combines lexical matching (e.g., BM25 or keyword filtering) with semantic similarity search to improve retrieval precision. Thakur et al. (2021) demonstrate that hybrid approaches consistently outperform purely dense retrieval across heterogeneous benchmarks. Reranking further refines candidate relevance using additional contextual signals. Cross-encoder rerankers achieve high accuracy but incur substantial computational overhead (Nogueira \& Cho, 2019); lightweight multi-factor reranking---combining similarity score, recency weighting, and keyword overlap---offers comparable grounding improvements at lower latency (Gao et al., 2023). In dynamic knowledge domains, recency-aware retrieval is particularly important to avoid outdated context injection.

\section{RAG Evaluation Methodologies}
\label{sec:lit_eval}

Evaluating RAG systems requires assessing both retrieval quality and generation faithfulness simultaneously. Es et al. (2023) introduced RAGAS, a reference-free evaluation framework that measures faithfulness, answer relevance, context precision, and context recall using automated LLM-based scoring, enabling scalable evaluation without large human annotation pools. Complementary work by Saad-Falcon et al. (2023) proposed ARES, which uses a small set of human-annotated examples to train lightweight judges capable of scoring RAG systems across multiple quality dimensions. These frameworks demonstrate that even compact evaluation sets---in the range of 20 to 50 annotated examples---can yield statistically meaningful comparisons between retrieval configurations when evaluation dimensions are clearly scoped and metrics are carefully defined (Es et al., 2023). This finding provides empirical support for the evaluation scale adopted in TechPulse.

\section{Data and Concept Drift in Retrieval Systems}
\label{sec:lit_drift}

In dynamic knowledge domains, retrieval system performance can degrade over time due to data drift---a shift in the statistical properties of incoming documents relative to the indexed corpus---or concept drift, where the semantic meaning of user queries evolves independently of corpus updates (Gama et al., 2014). Lu et al. (2018) provide a comprehensive taxonomy of drift detection methods, categorizing approaches into statistical process control, windowed distribution comparison, and ensemble-based disagreement detection. In retrieval-augmented systems specifically, drift can manifest as declining context precision when newly ingested documents diverge from the embedding space learned during the original indexing phase. Addressing drift requires both monitoring mechanisms that can detect distributional shifts in retrieval scores over time, and remediation actions such as corpus re-indexing, embedding model refresh, or query routing adjustments (Bayram et al., 2022). TechPulse incorporates a dual-criteria embedding distribution monitoring strategy informed by these findings.

\section{Hallucination Detection and Faithfulness Verification}
\label{sec:lit_hallucination}

Despite retrieval augmentation, generative models can still produce hallucinated content by inferring beyond the supplied context, over-generalizing retrieved passages, or introducing spurious cross-document associations (Ji et al., 2023). Manakul et al. (2023) proposed SelfCheckGPT, a consistency-based approach that detects hallucinations by sampling multiple outputs and measuring inter-sample consistency---a technique applicable without access to ground-truth references. Min et al. (2023) proposed FActScoring, which decomposes generated responses into atomic claims and individually verifies each claim against a reference corpus, offering fine-grained faithfulness assessment. At a system design level, Shi et al. (2023) demonstrate that irrelevant context injected alongside relevant passages can distract generative models and induce hallucinated responses even when correct information is present---motivating strict top-$k$ selection and high-precision retrieval as the primary anti-hallucination mechanism in TechPulse.

\section{MLOps and Production Machine Learning Systems}
\label{sec:lit_mlops}

Sculley et al. (2015) identify hidden technical debt as a central challenge in ML systems, emphasizing monitoring, reproducibility, and lifecycle management. MLOps extends DevOps principles to ML workflows, incorporating continuous integration, deployment automation, model versioning, and observability (Breck et al., 2017). Cloud-native, serverless architectures enable scalable and maintainable deployment with cost governance (Zhang et al., 2020; AWS, 2023). While RAG research frequently focuses on retrieval effectiveness, production-readiness---including pipeline idempotency, schema versioning, and budget-aware inference---remains underexplored.

\section{Research Gap}
\label{sec:gap}

The reviewed literature establishes six key insights: (1) RAG improves factual grounding but baseline dense retrieval may lack precision; (2) hybrid retrieval outperforms vector-only search in heterogeneous settings; (3) lightweight reranking can improve relevance without prohibitive overhead; (4) production ML systems require robust MLOps practices beyond model accuracy; (5) even compact evaluation sets of 20--50 annotated examples can yield statistically meaningful RAG comparisons when evaluation dimensions are clearly scoped; and (6) hallucination in RAG systems can persist despite retrieval augmentation, necessitating explicit faithfulness verification mechanisms. However, no existing work integrates all six dimensions into a unified, cost-aware, cloud-native system explicitly targeting dynamic emerging technology intelligence. This project addresses that gap.

Table~\ref{tab:related_work} summarizes representative related work and positions TechPulse relative to prior systems.

\begin{table}[H]
\caption{Comparison of Related Work and Positioning of TechPulse.}
\label{tab:related_work}
\centering
\resizebox{\textwidth}{!}{%
\begin{tabular}{p{2.8cm}p{1.8cm}p{1.3cm}p{1.6cm}p{1.6cm}p{1.3cm}p{1.5cm}p{1.8cm}}
\toprule
\textbf{System} & \textbf{Retrieval} & \textbf{Rerank} & \textbf{Real-Time} & \textbf{Cost Gov.} & \textbf{MLOps} & \textbf{Drift Det.} & \textbf{Halluc.\ Verif.} \\
\midrule
RAG (Lewis et al., 2020) & Dense & No & No & No & No & No & No \\
Atlas (Izacard et al., 2022) & Dense & No & No & No & No & No & No \\
RETRO (Borgeaud et al., 2022) & Dense & No & No & No & No & No & No \\
Typical RAG Prototype & Dense & Optional & Rare & No & Minimal & Rare & Rare \\
\textbf{TechPulse (ours)} & \textbf{Hybrid} & \textbf{Yes} & \textbf{Yes} & \textbf{Yes} & \textbf{Full} & \textbf{Dual} & \textbf{3-Layer} \\
\bottomrule
\end{tabular}%
}
\end{table}

\section{Chapter Summary}

This chapter reviewed the evolution of RAG architectures, embedding-based retrieval, hybrid retrieval optimization, MLOps principles, RAG evaluation methodologies, drift detection strategies, and hallucination verification approaches. The literature confirms that hybrid retrieval with lightweight reranking and production-grade MLOps integration represents an underexplored but practically important direction, and that compact but well-scoped evaluation sets can yield statistically meaningful RAG comparisons. The next chapter presents the methodology as implemented, including deviations from the original proposal.

% ============================================================
% CHAPTER 3: METHODOLOGY
% ============================================================
\chapter{Methodology}

\section{Overview of Technical Approach}
\label{sec:method_overview}

TechPulse has been implemented as a cloud-native Hybrid Retrieval-Augmented Generation system designed to synthesize emerging technology intelligence from heterogeneous sources. The technical approach is organized around five interconnected components: (1) continuous heterogeneous data ingestion from five sources; (2) deterministic preprocessing and semantic embedding; (3) hybrid multi-stage retrieval; (4) context-bound grounded generation with automatic LLM fallback; and (5) production-grade MLOps deployment and monitoring.

The core hypothesis under evaluation is that hybrid retrieval---combining structured metadata filtering, dense cosine similarity search, and lightweight multi-factor reranking---improves contextual precision and grounding faithfulness compared to baseline vector-only retrieval, while maintaining acceptable p95 latency and remaining within defined cost thresholds. This hypothesis is being tested through a controlled experiment comparing both configurations under identical corpus and embedding conditions.

\section{Why This Approach is Suitable}
\label{sec:method_rationale}

Three considerations motivate the hybrid RAG approach for this specific problem domain.

\textbf{Domain characteristics.} Emerging technology knowledge is temporally dynamic and distributed across heterogeneous sources. Vector-only retrieval treats all documents equally regardless of recency or source authority; hybrid retrieval with recency-aware reranking explicitly addresses temporal drift, which is critical in rapidly evolving AI/ML domains where a paper from 18 months ago may already be superseded.

\textbf{Operational constraints.} The system must operate within cost and latency bounds suitable for a research prototype deployed under AWS Academy Learner Lab restrictions. The lightweight linear reranking function (Equation~\ref{eq:reranking}) adds $O(k)$ computation over the top-$N$ filtered candidates---negligible compared to the $O(Nd)$ baseline search---and avoids the latency penalties of cross-encoder rerankers, which typically require a forward pass through a BERT-class model per candidate.

\textbf{Production feasibility.} The architecture uses AWS-managed services---Lambda, RDS PostgreSQL, API Gateway, EventBridge---eliminating the need to manage dedicated compute clusters or self-hosted vector databases. This enables a student team to deploy, iterate, and monitor a production-grade system within the course timeline and a prototype budget envelope.

\section{Technical Innovation and Optimization Strategy}
\label{sec:innovation}

The technical novelty of TechPulse lies not in any single component but in their integration into a coherent, cost-governed MLOps system. Specifically, four aspects represent the technical contribution of this project.

\textbf{Recency-aware hybrid reranking for dynamic knowledge domains.} The weighted linear reranking function (Equation~\ref{eq:reranking}) fuses semantic similarity, lexical keyword overlap, and exponential temporal decay into a single interpretable score. Unlike cross-encoder rerankers that are domain-agnostic and computationally expensive, this formulation is specifically designed for domains where content freshness is a first-class retrieval signal. The exponential decay parameter $\lambda$ provides a principled, tunable mechanism for controlling how rapidly older documents lose retrieval priority.

\textbf{Unified vector--relational retrieval without dedicated infrastructure.} By extending PostgreSQL with pgvector, TechPulse achieves hybrid retrieval (structured SQL metadata filtering + cosine vector search) within a single managed database, eliminating the operational complexity of running a separate vector store. This design choice reduces architectural surface area, simplifies deployment automation, and enables cost-governed scaling through standard RDS instances.

\textbf{Idempotent ingestion pipeline with SHA-256 content-addressed deduplication.} The SHA-256 hashing strategy applied to canonicalized document fields ensures that the pipeline can be re-run safely at any time without creating duplicate records. This is an often-overlooked but practically important data engineering innovation---it enables safe experiment replay, crash recovery, and corpus versioning without requiring a separate deduplication service.

\textbf{Budget-aware inference with programmatic halt.} The integration of AWS Cost Explorer queries into the inference path---where invocations are halted programmatically when monthly cost thresholds are approached---represents an operational optimization strategy not typically addressed in RAG research prototypes. This transforms cost from a passive monitoring concern into an active constraint enforced at runtime, consistent with production-grade MLOps principles.

\section{Baseline Retrieval Configuration}
\label{sec:baseline}

The baseline configuration performs dense cosine similarity search across the full indexed corpus. Given a query embedding $\mathbf{e}_q$ and document embeddings $\mathbf{e}_i$, cosine similarity is computed as:

\begin{equation}
  \text{Sim}(q, i) = \frac{\mathbf{e}_q \cdot \mathbf{e}_i}{\|\mathbf{e}_q\|\|\mathbf{e}_i\|}
  \label{eq:cosine}
\end{equation}

Top-$N$ documents are selected purely on similarity ranking with a minimum similarity threshold of 0.15 to filter retrieval noise. Under exact search, computational complexity scales as $O(Nd)$, where $N$ is the number of indexed chunks and $d = 384$ (the MiniLM embedding dimension). This configuration serves as the experimental control for retrieval performance comparison.

\section{Implemented Hybrid Retrieval Configuration}
\label{sec:hybrid}

The hybrid configuration introduces three sequential refinement stages applied before generation.

\textbf{Stage 1 --- Metadata Filtering.} Structured SQL constraints reduce the candidate space based on document recency (a 180-day sliding window), source type, and document state. Only documents in the \texttt{INDEXED} state are considered. This pre-filtering stage reduces retrieval noise prior to semantic ranking.

\textbf{Stage 2 --- Vector Similarity Search.} Cosine similarity (Equation~\ref{eq:cosine}) is applied over the filtered candidate subset using the pgvector \texttt{<=>} operator. To provide sufficient candidates for reranking, $4 \times \text{top\_k}$ results are retrieved at this stage.

\textbf{Stage 3 --- Reranking-Lite.} A lightweight weighted scoring function refines candidate ordering:

\begin{equation}
  \text{Score}_i = \alpha \cdot \text{Sim}_i + \beta \cdot \text{KeywordOverlap}_i + \gamma \cdot \text{RecencyWeight}_i
  \label{eq:reranking}
\end{equation}

where $\alpha + \beta + \gamma = 1$. The implemented default weights are $\alpha = 0.6$, $\beta = 0.2$, $\gamma = 0.2$, configurable via environment variables (\texttt{RERANK\_ALPHA}, \texttt{RERANK\_BETA}, \texttt{RERANK\_GAMMA}) to support grid search during evaluation.

Recency is modeled using exponential decay:

\begin{equation}
  \text{RecencyWeight}_i = e^{-\lambda \Delta t_i}
  \label{eq:recency}
\end{equation}

where $\Delta t_i$ is document age in days and $\lambda = 0.01$ (configurable via \texttt{RECENCY\_LAMBDA}).

Figure~\ref{fig:retrieval_comparison} illustrates the conceptual difference between the baseline and hybrid retrieval pathways.

% ---- TikZ: Baseline vs Hybrid comparison diagram ----
\begin{figure}[H]
\centering
\begin{tikzpicture}[node distance=0.45cm and 1.4cm]

% ---- LEFT: BASELINE ----
\node[pbox=awsgray, minimum width=2.6cm] (bq) {User Query};
\node[pbox=awsblue, below=of bq, minimum width=2.6cm] (bvs) {Vector Similarity\\Search (Full Corpus)};
\node[pbox=awsgreen, below=of bvs, minimum width=2.6cm] (btopk) {Top-$k$ Selection};
\node[pbox=awsorange, below=of btopk, minimum width=2.6cm] (bgen) {LLM\\Generation};

\draw[parrow] (bq) -- (bvs);
\draw[parrow] (bvs) -- (btopk);
\draw[parrow] (btopk) -- (bgen);

\node[layerlabel=awsgray, above=0.3cm of bq, font=\small\bfseries] {\textbf{Baseline (Vector-Only)}};

\begin{scope}[on background layer]
\node[boundbox=awsgray, fit=(bq)(bvs)(btopk)(bgen), inner sep=8pt] {};
\end{scope}

% ---- RIGHT: HYBRID ----
\node[pbox=awsgray, minimum width=2.6cm, right=3.6cm of bq] (hq) {User Query};
\node[pbox=propred, below=of hq, minimum width=2.6cm] (hmf) {Metadata Filter\\(recency + source)};
\node[pbox=awsblue, below=of hmf, minimum width=2.6cm] (hvs) {Vector Similarity\\Search (Filtered)};
\node[pbox=awspurple, below=of hvs, minimum width=2.6cm] (hrr) {Reranking-Lite\\($\alpha$\textperiodcentered Sim + $\beta$\textperiodcentered Kw + $\gamma$\textperiodcentered Rec)};
\node[pbox=awsgreen, below=of hrr, minimum width=2.6cm] (htopk) {Top-$k$ Selection};
\node[pbox=awsorange, below=of htopk, minimum width=2.6cm] (hgen) {LLM\\Generation};

\draw[parrow] (hq) -- (hmf);
\draw[parrow] (hmf) -- (hvs);
\draw[parrow] (hvs) -- (hrr);
\draw[parrow] (hrr) -- (htopk);
\draw[parrow] (htopk) -- (hgen);

\node[layerlabel=propblue, above=0.3cm of hq, font=\small\bfseries] {\textbf{Implemented Hybrid}};

\begin{scope}[on background layer]
\node[boundbox=propblue, fit=(hq)(hmf)(hvs)(hrr)(htopk)(hgen), inner sep=8pt] {};
\end{scope}

% annotation arrows
\draw[-{Stealth}, dashed, color=propred!70, thick]
  ([xshift=1.4cm]hmf.east) -- ++(0.5,0)
  node[right, font=\scriptsize, color=propred, text width=2.0cm]
  {Reduces candidate\\pool before search};

\draw[-{Stealth}, dashed, color=awspurple!70, thick]
  ([xshift=1.4cm]hrr.east) -- ++(0.5,0)
  node[right, font=\scriptsize, color=awspurple, text width=2.0cm]
  {Multi-factor\\relevance scoring};

\end{tikzpicture}
\caption{Comparison of Baseline Vector-Only Retrieval (left) and Implemented Hybrid Retrieval Pipeline (right). The hybrid pipeline adds a metadata filtering stage to reduce the candidate pool and a reranking-lite stage for multi-factor relevance scoring prior to final candidate selection.}
\label{fig:retrieval_comparison}
\end{figure}

\section{Hallucination Verification Strategy}
\label{sec:prompt}

TechPulse applies a three-layer verification approach to detect and flag potentially hallucinated outputs. The prompt construction and template details are presented in Section~\ref{sec:prompt_ch4}.

\begin{enumerate}[leftmargin=0.5in]
  \item \textbf{Prompt layer.} The system instruction enforces strict evidence grounding and abstention on insufficient context.

  \item \textbf{Evaluation layer.} RAGAS faithfulness scoring decomposes each generated response into atomic claims and verifies whether each claim is supported by at least one retrieved document.

  \item \textbf{Output layer --- Citation traceability.} At runtime, automated citation grounding verification analyzes each response sentence-by-sentence to compute a citation ratio---the fraction of sentences containing valid \texttt{[Source N]} references. Responses falling below the configurable citation grounding threshold are flagged as potentially hallucinated and replaced with an abstention message, with the original response preserved in metadata for review.
\end{enumerate}

\section{Optimization Objective}
\label{sec:optimization}

The retrieval task is framed as a constrained optimization problem:

\begin{equation}
  \max \; \text{Precision}(q) \quad \text{subject to} \quad \text{Latency} \leq \tau
  \label{eq:opt}
\end{equation}

where $\tau$ is the p95 latency target of 2 seconds. This formulation makes the precision--latency trade-off explicit and provides a principled basis for tuning reranking depth and metadata filter selectivity.

\section{Tool Selection and Implementation Stack}
\label{sec:tools}

Table~\ref{tab:tools} summarizes the implemented tools and their justification. The stack has been adapted from the original proposal to accommodate AWS Academy Learner Lab constraints, as detailed in Chapter~\ref{ch:progress}.

\begin{table}[H]
\caption{Implemented Tools, Services, and Justification.}
\label{tab:tools}
\centering
\begin{tabular}{p{3.2cm}p{3.2cm}p{6.6cm}}
\toprule
\textbf{Component} & \textbf{Tool / Service} & \textbf{Justification} \\
\midrule
Ingestion orchestration & AWS Lambda + EventBridge & Serverless, scheduled/event-driven, no cluster management \\
Message queue & Amazon SQS + DLQ & Decouples ingestion from embedding; failure resilience \\
Raw storage & Amazon S3 (medallion layout) & Cost-effective, durable, prefix-based lineage tracking \\
Embedding model & MiniLM (sentence-transformers) & Low latency, 384-dim, serverless-compatible \\
Vector + relational store & RDS PostgreSQL + pgvector & Unified hybrid retrieval; Learner Lab compatible \\
Generative model & Groq (default) / Bedrock / HuggingFace / Ollama & Multi-backend auto-cascade; Groq free tier primary \\
Frontend & React (S3 static hosting) & Lightweight, decoupled from backend \\
API routing & Amazon API Gateway (HTTP) & Managed, HTTPS, CORS support \\
CI/CD & GitHub Actions + SAM & Reproducible IaC deployments, automated testing \\
Monitoring & CloudWatch + SNS alarms & Observability, error alerts, health checks \\
Evaluation & RAGAS + grid search & Standardized RAG quality metrics + weight optimization \\
\bottomrule
\end{tabular}
\end{table}

% ============================================================
% CHAPTER 4: SYSTEM ARCHITECTURE AND PIPELINE DESIGN
% ============================================================
\chapter{System Architecture and Pipeline Design}

\section{Architectural Overview}
\label{sec:arch_overview}

TechPulse has been implemented as a six-layer, cloud-native system adapted for deployment within the AWS Academy Learner Lab environment. Each layer is logically independent with well-defined interfaces, enabling modular development, testing, and future extension. The architecture has been adapted from the original proposal to replace services unavailable in Learner Lab (Aurora Serverless, Amazon Bedrock, custom VPCs, AWS Budgets) with compatible alternatives while preserving the core hybrid retrieval and MLOps capabilities.

% ---- TikZ: Six-layer architectural stack ----
\begin{figure}[H]
\centering
\begin{tikzpicture}[node distance=0.0cm]

\def\lw{13.5cm}
\def\lh{0.90cm}

\node[rectangle, draw=awsgray!50, fill=awslightblue, minimum width=\lw,
      minimum height=\lh, font=\small\bfseries, text centered,
      rounded corners=2pt]
  (l6) {Layer 6 --- Deployment, Monitoring \& Interaction \quad \textcolor{awsgray}{\scriptsize CloudWatch, SNS, API GW, React}};

\node[rectangle, draw=awsgray!50, fill=awslightorange, minimum width=\lw,
      minimum height=\lh, font=\small\bfseries, text centered,
      rounded corners=2pt, below=0pt of l6]
  (l5) {Layer 5 --- Generative Synthesis (Groq / Bedrock / Ollama / HuggingFace) \quad \textcolor{awsgray}{\scriptsize Auto-fallback chain}};

\node[rectangle, draw=awsgray!50, fill=awslightpurple, minimum width=\lw,
      minimum height=\lh, font=\small\bfseries, text centered,
      rounded corners=2pt, below=0pt of l5]
  (l4) {Layer 4 --- Hybrid Retrieval \& Reranking \quad \textcolor{awsgray}{\scriptsize 3-stage: filter $\to$ similarity $\to$ rerank}};

\node[rectangle, draw=awsgray!50, fill=awslightgreen, minimum width=\lw,
      minimum height=\lh, font=\small\bfseries, text centered,
      rounded corners=2pt, below=0pt of l4]
  (l3) {Layer 3 --- Vector Storage \& Indexing (RDS PostgreSQL + pgvector) \quad \textcolor{awsgray}{\scriptsize HNSW index, cosine search + SQL filter}};

\node[rectangle, draw=awsgray!50, fill=awslightorange, minimum width=\lw,
      minimum height=\lh, font=\small\bfseries, text centered,
      rounded corners=2pt, below=0pt of l3]
  (l2) {Layer 2 --- Preprocessing \& Feature Engineering \quad \textcolor{awsgray}{\scriptsize 17-step normalize, chunk, SHA-256 hash}};

\node[rectangle, draw=awsgray!50, fill=awslightblue, minimum width=\lw,
      minimum height=\lh, font=\small\bfseries, text centered,
      rounded corners=2pt, below=0pt of l2]
  (l1) {Layer 1 --- Data Ingestion (5 sources) \quad \textcolor{awsgray}{\scriptsize ArXiv + HN + DEV.to + GitHub + RSS (batch + near-RT)}};

\end{tikzpicture}
\caption{Six-Layer Conceptual Architecture of TechPulse as implemented. Layers run from data ingestion (bottom) to user interaction and monitoring (top), each with well-defined responsibilities.}
\label{fig:six_layers}
\end{figure}

\section{Alignment with the AWS Well-Architected Framework}
\label{sec:well_architected}

The TechPulse architecture is designed in explicit alignment with the AWS Well-Architected Framework (AWS, 2023). Table~\ref{tab:well_arch} summarises how each pillar is addressed within the implemented system.

\begin{table}[H]
\caption{TechPulse Alignment with the AWS Well-Architected Framework Six Pillars.}
\label{tab:well_arch}
\centering
\begin{tabular}{p{3.2cm}p{9.3cm}}
\toprule
\textbf{Pillar} & \textbf{TechPulse Implementation} \\
\midrule
\textbf{Operational Excellence} & Automated ingestion via EventBridge (6-hour cycle); CI/CD pipeline using GitHub Actions and SAM; CloudWatch logging, custom metrics, and health-check Lambda (5-minute interval) \\[4pt]
\textbf{Security} & Controlled API access via Amazon API Gateway with CORS; rate limiting via slowapi; all services managed within AWS default VPC; no direct model invocation from the frontend \\[4pt]
\textbf{Reliability} & Decoupled ingestion and embedding pipelines via SQS; Dead Letter Queue for failure recovery and replay with 14-day retention; bounded retry logic with exponential backoff for LLM invocations; automatic LLM backend fallback chain \\[4pt]
\textbf{Performance Efficiency} & Serverless Lambda scaling; lightweight MiniLM embedding model (384-dim); metadata pre-filtering reduces vector search scope; HNSW index for approximate nearest-neighbor search \\[4pt]
\textbf{Cost Optimization} & Serverless-only execution (pay-per-invocation); programmatic budget guard via Cost Explorer; Groq free-tier inference as primary backend in Learner Lab; db.t3.micro RDS instance \\[4pt]
\textbf{Sustainability} & Event-driven processing avoids idle compute; SQS decoupling prevents over-provisioning; right-sized Lambda memory allocations (512MB--1024MB) \\
\bottomrule
\end{tabular}
\end{table}

\section{Data Sources}
\label{sec:data_sources}

The original proposal specified two data sources (ArXiv and Hacker News). During implementation, the data source coverage was expanded to five heterogeneous sources to improve corpus diversity, topical breadth, and temporal coverage. Table~\ref{tab:data_sources} summarizes the implemented data sources.

\begin{table}[H]
\caption{Implemented Data Sources and Roles.}
\label{tab:data_sources}
\centering
\begin{tabular}{p{3.0cm}p{2.5cm}p{4.5cm}p{2.5cm}}
\toprule
\textbf{Source} & \textbf{Type} & \textbf{Role} & \textbf{Frequency} \\
\midrule
ArXiv API & Batch & Scholarly grounding corpus (cs.AI, cs.LG, cs.CL, cs.IR, cs.SE) & Every 12 hours \\
Hacker News API & Near-real-time & Community-curated tech discussions & Every 30 minutes \\
DEV.to API & Near-real-time & Practitioner tutorials and industry articles (21 tags) & Every 6 hours \\
GitHub Trending & Batch & Open-source project intelligence (10 topic queries) & Every 6 hours \\
RSS/Atom Feeds & Batch & Curated technology journalism (10 feeds) & Every 6 hours \\
\bottomrule
\end{tabular}
\end{table}

The combination provides complementary temporal and authority profiles: ArXiv contributes rigorously authored research with stable identifiers; Hacker News contributes practitioner signals and community-curated trending discussions; DEV.to provides developer-focused tutorials and implementation insights; GitHub Trending captures emerging open-source tools and frameworks; and curated RSS feeds from outlets including Ars Technica, TechCrunch, Wired, IEEE Spectrum, and MIT News provide technology journalism coverage.

\section{Data Engineering Pipeline}
\label{sec:pipeline}

The pipeline has been implemented around four design principles: determinism (identical input produces identical output), idempotency (re-running does not create duplicate records), observability (all state transitions are logged), and failure resilience (failures are isolated and recoverable).

% ---- TikZ: Full data engineering pipeline ----
\begin{figure}[H]
\centering
\resizebox{\textwidth}{!}{%
\begin{tikzpicture}[node distance=0.4cm and 0.5cm, font=\small]

% ---- Sources ----
\node[compnode=awsblue] (arxiv) {ArXiv API\\(12 h)};
\node[compnode=awsblue, below=0.25cm of arxiv] (hn) {Hacker News\\(30 min)};
\node[compnode=awsblue, below=0.25cm of hn] (devto) {DEV.to\\(6 h)};
\node[compnode=awsblue, below=0.25cm of devto] (github) {GitHub\\(6 h)};
\node[compnode=awsblue, below=0.25cm of github] (rss) {RSS Feeds\\(6 h)};

% ---- EventBridge ----
\node[compnode=awsorange, right=0.8cm of hn, yshift=-0.55cm] (eb) {EventBridge\\Scheduler};

% ---- Ingestion Lambda ----
\node[compnode=awsorange, right=0.7cm of eb] (inlambda) {Ingestion\\Lambda};

% ---- S3 Raw ----
\node[compnode=awsgreen, above right=0.1cm and 0.7cm of inlambda] (s3raw) {S3 Raw\\(5 source prefixes)};

% ---- SQS ----
\node[compnode=awsgray, right=0.7cm of inlambda, yshift=-0.55cm] (sqs) {SQS\\Queue};

% ---- DLQ ----
\node[compnode=propred!80, below=0.35cm of sqs] (dlq) {Dead Letter\\Queue};

% ---- Preprocess + Embed Lambda ----
\node[compnode=awsorange, right=0.7cm of sqs] (pelambda) {Preprocess \&\\Embed Lambda\\(MiniLM)};

% ---- S3 Processed ----
\node[compnode=awsgreen, above right=0.1cm and 0.7cm of pelambda] (s3proc) {S3 Processed\\(/processed)};

% ---- RDS ----
\node[compnode=awspurple, right=0.7cm of pelambda, yshift=-0.55cm] (rds) {RDS PG\\+ pgvector};

% ---- Retrieval / LLM ----
\node[compnode=awsorange, right=0.7cm of rds] (rag) {RAG\\Orchestrator};
\node[compnode=awsblue, right=0.7cm of rag] (llm) {LLM Backend\\(Groq/Bedrock)};

% ---- API GW + React ----
\node[compnode=awsgreen, right=0.7cm of llm] (apigw) {API\\Gateway};
\node[compnode=awsblue, right=0.5cm of apigw] (react) {React\\Frontend};

% ---- Observability ----
\node[compnode=propgray, below=0.8cm of rag, xshift=0.5cm] (cw) {CloudWatch\\Metrics/Logs};
\node[compnode=propgray, right=0.6cm of cw] (sns) {SNS\\Alerts};
\node[compnode=propgray, left=0.6cm of cw] (hc) {Health Check\\Lambda};

% Arrows: sources -> eb -> inlambda
\draw[parrow] (arxiv.east) -- ++(0.25,0) |- (eb.west);
\draw[parrow] (hn.east)    -- ++(0.25,0) |- (eb.west);
\draw[parrow] (devto.east) -- ++(0.25,0) |- (eb.west);
\draw[parrow] (github.east) -- ++(0.25,0) |- (eb.west);
\draw[parrow] (rss.east)   -- ++(0.25,0) |- (eb.west);
\draw[parrow] (eb) -- (inlambda);
\draw[parrow] (inlambda.east) -- ++(0.25,0) |- (s3raw.west);
\draw[parrow] (inlambda) -- (sqs);
\draw[parrow, dashed, color=propred] (sqs.south) -- (dlq.north)
  node[midway, right, font=\scriptsize, color=propred] {failures};
\draw[parrow] (sqs) -- (pelambda);
\draw[parrow] (pelambda.east) -- ++(0.25,0) |- (s3proc.west);
\draw[parrow] (pelambda) -- (rds);
\draw[parrow] (rds) -- (rag);
\draw[parrow] (rag) -- (llm);
\draw[parrow] (llm) -- (apigw);
\draw[parrow] (apigw) -- (react);

% Observability arrows
\draw[parrow, dashed, color=propgray] (rag.south) -- (cw.north);
\draw[parrow, dashed, color=propgray] (cw.east) -- (sns.west);
\draw[parrow, dashed, color=propgray] (hc.east) -- (cw.west);

% Layer labels
\node[below=0.5cm of devto, font=\scriptsize\bfseries, color=awsblue, xshift=0.2cm]
  {Sources};
\node[below=0.5cm of eb, font=\scriptsize\bfseries, color=awsorange]
  {Trigger};
\node[below=1.2cm of sqs, font=\scriptsize\bfseries, color=awsorange, xshift=0.8cm]
  {Processing};
\node[below=0.5cm of rds, font=\scriptsize\bfseries, color=awspurple]
  {Storage};
\node[below=0.5cm of rag, font=\scriptsize\bfseries, color=awsorange, xshift=0.3cm]
  {Serving};
\node[below=0.3cm of cw, font=\scriptsize\bfseries, color=propgray]
  {Observability};

\end{tikzpicture}%
}
\caption{Implemented End-to-End Data Engineering and Serving Pipeline for TechPulse. Five heterogeneous sources feed into the ingestion pipeline. Solid arrows indicate data and control flow; dashed arrows indicate monitoring and observability signals.}
\label{fig:pipeline}
\end{figure}

\subsection{Ingestion and Decoupling}

Ingestion has been implemented through AWS Lambda functions orchestrated by Amazon EventBridge, supporting both scheduled and event-driven execution. Each of the five ingesters (ArXiv, Hacker News, DEV.to, GitHub, RSS) operates independently with source-specific data extraction logic. Raw artifacts are persisted in Amazon S3 under a structured prefix hierarchy (\texttt{/raw/arxiv/}, \texttt{/raw/hn/}, \texttt{/raw/devto/}, \texttt{/raw/github/}, \texttt{/raw/rss/}, \texttt{/processed/}) to preserve transformation lineage. The ingestion pipeline is decoupled from the embedding pipeline via Amazon SQS, providing backpressure control, failure isolation, and independent horizontal scaling. Messages that fail processing after three retry attempts are routed to a Dead Letter Queue with 14-day retention for controlled replay.

All five ingesters share a common architectural pattern: (1) fetch data from API or feed; (2) extract and normalize fields; (3) compute SHA-256 content hash over canonicalized fields; (4) insert into the \texttt{documents} table with \texttt{ON CONFLICT (content\_hash) DO NOTHING} for idempotent deduplication; (5) write raw JSON to S3; and (6) enqueue document ID to SQS for downstream processing.

\subsection{Document State Lifecycle}

Each document progresses through a four-state deterministic lifecycle to guarantee idempotency and traceability. Table~\ref{tab:doc_states} defines these states.

\begin{table}[H]
\caption{Document State Transition Model.}
\label{tab:doc_states}
\centering
\begin{tabular}{p{2.0cm}p{3.8cm}p{3.0cm}p{3.0cm}}
\toprule
\textbf{State} & \textbf{Description} & \textbf{Storage} & \textbf{Trigger} \\
\midrule
RAW & Document ingested from API & Amazon S3 & Successful API fetch \\
PROCESSED & Cleaned, normalized, chunked & S3 (\texttt{/processed}) & Preprocessing complete \\
EMBEDDED & Vector representation generated & RDS (pgvector) & Embedding computed \\
INDEXED & Available for retrieval & RDS metadata table & Vector committed \\
\bottomrule
\end{tabular}
\end{table}

To guarantee idempotency, canonicalized content fields are hashed using SHA-256 before processing. The hash covers source-specific canonical fields: for ArXiv, title, abstract, and publication date; for Hacker News, title, URL, and timestamp; for DEV.to, article title, URL, and tags; for GitHub, repository full name, description, and stars; for RSS, entry title, link, and published date. Duplicate hashes are skipped at ingestion via PostgreSQL \texttt{ON CONFLICT} clauses, ensuring consistent corpus integrity across experimental runs.

\subsection{Preprocessing and Chunking}

Textual preprocessing has been implemented as a comprehensive 17-step normalization pipeline in the \texttt{normalize\_text()} function. The following transformations are applied in sequence:

\begin{enumerate}[leftmargin=0.5in]
  \item HTML entity unescaping (\texttt{html.unescape})
  \item HTML comment removal (\texttt{<!-- ... -->})
  \item HTML tag stripping
  \item Fenced code block removal
  \item Inline code formatting removal
  \item Markdown image and link syntax normalization
  \item URL removal (HTTP/HTTPS/FTP)
  \item Markdown heading marker removal
  \item Bold and italic marker removal
  \item Strikethrough syntax normalization (\texttt{\~{}\~{}text\~{}\~{}} $\to$ text)
  \item Horizontal rule removal
  \item Blockquote marker removal
  \item Emoji and pictographic symbol removal (6 Unicode ranges)
  \item Unicode control character and zero-width character removal
  \item Repeated punctuation collapse (\texttt{!!!} $\to$ \texttt{!})
  \item Non-breaking space normalization
  \item Whitespace consolidation and trimming
\end{enumerate}

Following normalization, documents are segmented using token-based chunking with the \texttt{tiktoken} tokenizer (\texttt{cl100k\_base} encoding). The implemented configuration uses a 500-token window with 50-token overlap, balancing semantic coherence with retrieval granularity.

\section{Embedding and Indexing}
\label{sec:embedding}

Semantic representations are generated using the \texttt{all-MiniLM-L6-v2} sentence-transformer model (Wang et al., 2020), producing 384-dimensional embeddings. The model is loaded as a lazy-initialized singleton to minimize memory overhead under serverless execution constraints. Embeddings are stored in RDS PostgreSQL using the pgvector extension with an HNSW (Hierarchical Navigable Small World) index configured with $m = 16$ and $\text{ef\_construction} = 64$, enabling efficient approximate nearest-neighbor search with cosine distance.

\section{Storage Architecture (Medallion Layout)}
\label{sec:storage}

\begin{table}[H]
\caption{Storage Layers and Responsibilities (Medallion Architecture).}
\label{tab:storage_layers}
\centering
\begin{tabular}{p{2.8cm}p{5.8cm}p{4.0cm}}
\toprule
\textbf{Layer} & \textbf{Technology} & \textbf{Purpose} \\
\midrule
Raw (Bronze) & Amazon S3 (\texttt{/raw/\{source\}/}) & Store unprocessed API responses \\
Processed (Silver) & Amazon S3 (\texttt{/processed/}) & Store cleaned and chunked documents \\
Vector (Gold) & RDS PostgreSQL + pgvector & Store embeddings and structured metadata \\
\bottomrule
\end{tabular}
\end{table}

\section{Database Schema}
\label{sec:schema}

The PostgreSQL database schema implements three core tables. The \texttt{documents} table stores ingested records with fields for source, title, content, URL, publication date, SHA-256 content hash (unique constraint), document state, and timestamps. The \texttt{chunks} table stores preprocessed text segments with their 384-dimensional vector embeddings, linked to parent documents via foreign key. The \texttt{drift\_baselines} table persists drift detection history including mean similarity, standard deviation, probe count, and alert status.

Four database indexes support efficient retrieval: an HNSW vector index on chunk embeddings (\texttt{vector\_cosine\_ops}, $m=16$, $\text{ef\_construction}=64$), and B-tree indexes on \texttt{published\_at} (descending), \texttt{source}, and \texttt{state} for metadata filtering performance.

\section{Prompt Construction and Grounded Generation}
\label{sec:prompt_ch4}

Retrieved top-$k$ documents are concatenated into a structured prompt template with four logical sections: (1) \textbf{System Instruction}---defines assistant behavior, enforces evidence-grounded responses, and mandates \texttt{[Source N]} citation format; (2) \textbf{Context Block}---concatenates retrieved document excerpts with source identifiers, titles, and publication dates; (3) \textbf{User Question}---the original user query; and (4) \textbf{Citation Instruction}---requires the model to reference supporting documents.

\begin{lstlisting}[caption={Implemented Structured Generation Prompt Template.}, label={lst:prompt}]
SYSTEM:
Answer using the provided context documents as your
primary evidence. For each specific claim, cite the
source in [Source N] format. Do NOT speculate beyond
what the documents support.

CONTEXT:
[Source 1] Title: ... | Date: ... | Source: arxiv
Excerpt text...

[Source 2] Title: ... | Date: ... | Source: hn
Excerpt text...

QUESTION:
{User Query}

INSTRUCTION:
Provide a concise answer and cite supporting sources.
\end{lstlisting}

If retrieval returns insufficient evidence or if the citation grounding check fails, the system returns an abstention response rather than generating unsupported content.

\section{LLM Backend Fallback Chain}
\label{sec:llm_fallback}

To ensure generation availability across diverse deployment environments, TechPulse implements an automatic LLM backend fallback chain. When the primary backend fails (e.g., due to service unavailability, rate limiting, or network errors), the system automatically cascades to the next available backend without manual intervention.

\begin{table}[H]
\caption{Implemented LLM Backend Fallback Hierarchy.}
\label{tab:llm_backends}
\centering
\begin{tabular}{p{0.5cm}p{3.5cm}p{4.0cm}p{4.0cm}}
\toprule
\textbf{Tier} & \textbf{Backend} & \textbf{Fallback Chain} & \textbf{Notes} \\
\midrule
1 & Groq (llama-3.1-8b-instant) & $\to$ Bedrock $\to$ Ollama $\to$ HuggingFace & Primary backend; free tier, OpenAI-compatible API \\
2 & Amazon Bedrock (Converse API) & $\to$ Ollama $\to$ HuggingFace & AWS-native; model-agnostic via Converse API \\
3 & Ollama (local) & $\to$ HuggingFace & Local inference for development \\
4 & HuggingFace Inference API & $\to$ Ollama & Mistral-7B-Instruct; serverless \\
\bottomrule
\end{tabular}
\end{table}

All four backends produce responses compatible with the same structured prompt template, ensuring that evaluation results remain comparable across backends. Each backend is configured with bounded retry logic (2 retries with exponential backoff) before cascading to the next tier. If all backends fail, the system degrades gracefully by returning a retrieval-only summary constructed from the top-$k$ documents without generative expansion.

\section{Deployment, Monitoring, and MLOps}
\label{sec:mlops}

Figure~\ref{fig:mlops_loop} presents the implemented MLOps lifecycle loop.

% ---- TikZ: MLOps lifecycle loop ----
\begin{figure}[H]
\centering
\resizebox{0.92\textwidth}{!}{%
\begin{tikzpicture}[node distance=1.4cm and 0.5cm]

% Main cycle nodes
\node[pbox=awsblue, minimum width=2.8cm] (ingest) {Data Ingestion\\(EventBridge)};
\node[pbox=awsgreen, right=2.0cm of ingest] (process) {Preprocess \&\\Embed};
\node[pbox=awspurple, right=2.0cm of process] (index) {Index\\(pgvector)};
\node[pbox=awsorange, below=1.8cm of index] (serve) {Serve\\(LLM + RAG)};
\node[pbox=propblue, below=1.8cm of ingest] (monitor) {Monitor\\(CloudWatch)};
\node[pbox=propgreen, below=1.8cm of process] (evaluate) {Evaluate\\(RAGAS + Grid)};

% Cycle arrows
\draw[parrow, thick] (ingest) -- (process);
\draw[parrow, thick] (process) -- (index);
\draw[parrow, thick] (index) -- (serve);
\draw[parrow, thick] (serve) -- (evaluate);
\draw[parrow, thick] (evaluate) -- (monitor);
\draw[parrow, thick] (monitor) -- (ingest);

% CI/CD overlay
\node[pbox=propgray, above=1.1cm of process, minimum width=5.0cm] (cicd)
  {CI/CD: GitHub Actions + SAM};
\draw[-{Stealth}, dashed, color=propgray, thick] (cicd.south) -- (process.north);
\draw[-{Stealth}, dashed, color=propgray, thick] (cicd.south) -- ++(0,-0.3) -| (index.north);

% Budget guard
\node[pbox=propred, below=0.8cm of serve, minimum width=2.6cm] (budget)
  {Budget Guard\\(Cost Explorer)};
\draw[-{Stealth}, dashed, color=propred, thick] (serve.south) -- (budget.north);
\draw[-{Stealth}, dashed, color=propred, thick] (budget.south) -- ++(0,-0.5) -| (monitor.south);

\end{tikzpicture}%
}
\caption{Implemented MLOps lifecycle loop: continuous ingestion, preprocessing, indexing, serving, evaluation, and monitoring, with CI/CD automation and a budget governance guard.}
\label{fig:mlops_loop}
\end{figure}

\subsection{LLM Generation Configuration}

Generation is performed through the configured LLM backend with the following parameter bounds:

\begin{table}[H]
\caption{Implemented LLM Generation Parameter Configuration.}
\label{tab:llm_params}
\centering
\begin{tabular}{lll}
\toprule
\textbf{Parameter} & \textbf{Value} & \textbf{Rationale} \\
\midrule
Max Output Tokens & 300 & Cost control \\
Temperature & 0.2 & Reduce hallucination risk \\
Timeout & 30 seconds (Groq) / 300 seconds (Ollama) / 60 seconds (HF) & Backend-appropriate limits \\
Retry Attempts & 2 & Transient failure handling \\
Backoff Base & 2 seconds & Exponential backoff \\
Context Window & 4096 tokens (Ollama) & Model-appropriate context \\
\bottomrule
\end{tabular}
\end{table}

If all LLM backends fail after exhausting retries and fallback chains, the system degrades gracefully by returning a retrieval-only summary and logging the error to CloudWatch.

\subsection{Drift Detection and Remediation}
\label{sec:drift}

TechPulse implements a dual-criteria drift detection mechanism to identify degradation in retrieval quality over time:

\textbf{Probe query baseline.} At each drift detection cycle, 20 fixed canonical probe queries (stored in \texttt{probe\_queries.json}) are executed through the hybrid retrieval pipeline. The mean and standard deviation of cosine similarity scores across all probes are computed and persisted in the \texttt{drift\_baselines} database table.

\textbf{Detection criterion 1 --- Relative drop threshold.} The current mean similarity is compared against the most recent non-alert baseline. A drift alert is triggered if the relative drop exceeds 10\%.

\textbf{Detection criterion 2 --- Shewhart $3\sigma$ control chart.} The system retrieves the last 10 non-alert baselines via \texttt{\_get\_baseline\_history(n=10)} and computes the historical mean ($\mu$) and standard deviation ($\sigma$). The lower control limit is defined as:

\begin{equation}
  \text{LCL} = \mu_{\text{hist}} - 3\sigma_{\text{hist}}
  \label{eq:shewhart}
\end{equation}

A drift alert is triggered if the current mean similarity falls below LCL. This statistical process control approach detects systematic degradation that may not be captured by point-wise comparison alone.

Both criteria are evaluated independently, and either triggering constitutes a drift alert. The \texttt{alert\_reason} field records a human-readable explanation (e.g., ``relative drop 15.2\%; Shewhart $3\sigma$ breach (LCL=0.4321, value=0.3890)'') routed through CloudWatch metrics and SNS notification.

\textbf{Remediation protocol.} The ordered response upon drift alert includes: (1) corpus audit to identify localized drift by source or topic cluster; (2) metadata filter re-calibration if off-topic documents enter the corpus; (3) re-indexing affected corpus segments with updated preprocessing; and (4) embedding model refresh (deferred to post-prototype phase).

\subsection{Health Monitoring}

A scheduled health-check Lambda runs every 5 minutes via EventBridge, probing four service dependencies:

\begin{enumerate}[leftmargin=0.5in]
  \item \textbf{Database connectivity} --- executes a lightweight test query against RDS PostgreSQL and reports document and chunk counts.
  \item \textbf{S3 reachability} --- verifies the data bucket is accessible.
  \item \textbf{SQS reachability} --- confirms the ingestion queue is operational.
  \item \textbf{LLM backend reachability} --- sends a minimal test prompt (``Respond with OK.'') to the configured LLM backend and reports which backend responded.
\end{enumerate}

Results are logged to CloudWatch. Three CloudWatch alarms are configured: RAG API errors ($>3$ in 5 minutes), Ingestion errors ($>2$ in 5 minutes), and DLQ depth ($\geq 1$ message), all routing to SNS for operational notification.

\subsection{CI/CD and Deployment Automation}

A four-stage CI/CD pipeline has been implemented using GitHub Actions:

\begin{enumerate}[leftmargin=0.5in]
  \item \textbf{Lint and Test} --- Python 3.11 setup, dependency installation, Ruff linting across \texttt{src/} and \texttt{tests/}, Pytest execution with coverage reporting.
  \item \textbf{Docker Build} --- Builds \texttt{techpulse-backend} and \texttt{techpulse-frontend} container images.
  \item \textbf{SAM Build} --- Builds and validates the Learner Lab SAM template (\texttt{template-learnerlab.yaml}).
  \item \textbf{Deploy (main branch only)} --- Deploys to AWS via SAM with parameterized database credentials, VPC configuration, and LabRole ARN from GitHub Secrets.
\end{enumerate}

All infrastructure is defined as code using AWS SAM, ensuring reproducible and consistent deployments. The SAM template defines: RDS PostgreSQL instance (\texttt{db.t3.micro}), S3 data and frontend buckets, SQS queue with DLQ, four Lambda functions (Ingestion, Preprocess, RAG API, Health Check), HTTP API Gateway, SNS alert topic, EventBridge schedules, and CloudWatch alarms.

\section{User Interface}
\label{sec:ui}

The user interface is implemented as a lightweight React single-page application. It provides natural language query submission, structured response rendering with citation display, and retrieval metadata transparency. All responses are constrained through the RAG pipeline; direct model invocation is not permitted. When retrieval evidence is insufficient or hallucination is detected, the system returns an explicit fail-safe abstention message.

% ============================================================
% CHAPTER 5: EVALUATION FRAMEWORK
% ============================================================
\chapter{Experimental Plan and Evaluation Metrics}

\section{Evaluation Objective}
\label{sec:eval_objective}

The evaluation aims to answer three questions: (1) Does hybrid retrieval improve contextual precision and grounding faithfulness compared to vector-only retrieval? (2) What latency and cost overhead does hybrid retrieval introduce? (3) Does the system remain within acceptable operational thresholds under moderate workload?

\section{Experimental Configurations}
\label{sec:eval_configs}

Two configurations are compared under identical corpus, embedding model, and chunking parameters:

\begin{table}[H]
\caption{Experimental Configuration Comparison.}
\label{tab:exp_configs}
\centering
\begin{tabular}{p{3.5cm}cc}
\toprule
\textbf{Component} & \textbf{Baseline (Vector-Only)} & \textbf{Hybrid (Implemented)} \\
\midrule
Metadata Filtering & $\times$ & \checkmark \\
Vector Similarity & \checkmark & \checkmark \\
Reranking-Lite & $\times$ & \checkmark \\
Recency Weighting & $\times$ & \checkmark \\
Top-$k$ & 8 & 8 \\
Embedding Model & MiniLM (384-dim) & MiniLM (384-dim) \\
Cost Governance & None & Budget guard \\
\bottomrule
\end{tabular}
\end{table}

\section{Evaluation Dataset}
\label{sec:eval_dataset}

The evaluation framework uses 25 manually curated trend-oriented queries spanning recent research developments, emerging tools, and comparative technology discussions in artificial intelligence and related domains. Queries are categorized into two types---\textit{research\_trend} and \textit{tool\_technology}---and include ground-truth reference answers for automated scoring.

Example queries include: ``What are recent approaches to efficient transformer architectures?''; ``What emerging AI tools are currently being discussed by developers?''; and ``What are recent trends in large language model compression?''

\textbf{Justification for evaluation set size.} The evaluation objective is a controlled comparison between two retrieval configurations under identical corpus and embedding conditions. A paired comparison across 25 carefully constructed queries with ground-truth annotations is sufficient to detect meaningful precision differences, consistent with the evaluation scale adopted in prior RAG evaluation work (Es et al., 2023; Saad-Falcon et al., 2023). Queries are stratified across query categories to ensure results are not driven by a single query type.

\section{Evaluation Metrics}
\label{sec:eval_metrics}

Figure~\ref{fig:eval_framework} summarizes the implemented evaluation framework.

% ---- TikZ: Evaluation framework diagram ----
\begin{figure}[H]
\centering
\resizebox{\textwidth}{!}{%
\begin{tikzpicture}

\node[evalbox=awsblue] (prec) at (-6.0,0)
  {\textbf{Context Precision}\\[2pt]\scriptsize RAGAS automated\\scoring\\Target: $> 0.75$};

\node[evalbox=awsgreen] (faith) at (-2.0,0)
  {\textbf{Faithfulness}\\[2pt]\scriptsize RAGAS atomic claim\\verification\\High target};

\node[evalbox=awsorange] (lat) at (2.0,0)
  {\textbf{p95 Latency}\\[2pt]\scriptsize End-to-end + per-stage\\decomposition\\$< 2$ sec};

\node[evalbox=awspurple] (cost) at (6.0,0)
  {\textbf{Cost / Query}\\[2pt]\scriptsize Token usage + Lambda\\execution\\$<$USD 15/mo};

\node[rectangle, draw=awsgray!40, fill=sectionbg, rounded corners=3pt,
      font=\scriptsize, minimum height=0.55cm, text width=3.2cm, text centered]
  (m1) at (-6.0,-1.7) {RAGAS context precision\\without reference};

\node[rectangle, draw=awsgray!40, fill=sectionbg, rounded corners=3pt,
      font=\scriptsize, minimum height=0.55cm, text width=3.2cm, text centered]
  (m2) at (-2.0,-1.7) {Citation grounding\\ratio + RAGAS faithfulness};

\node[rectangle, draw=awsgray!40, fill=sectionbg, rounded corners=3pt,
      font=\scriptsize, minimum height=0.55cm, text width=3.2cm, text centered]
  (m3) at (2.0,-1.7) {Per-query timing\\+ CloudWatch Logs};

\node[rectangle, draw=awsgray!40, fill=sectionbg, rounded corners=3pt,
      font=\scriptsize, minimum height=0.55cm, text width=3.2cm, text centered]
  (m4) at (6.0,-1.7) {Token count + \\estimated cost};

\draw[parrow] (prec) -- (m1);
\draw[parrow] (faith) -- (m2);
\draw[parrow] (lat) -- (m3);
\draw[parrow] (cost) -- (m4);

\end{tikzpicture}%
}
\caption{Implemented evaluation framework: four dimensions---context precision, faithfulness, p95 latency, and cost per query---each with a measurement method and target threshold.}
\label{fig:eval_framework}
\end{figure}

\subsection{Context Precision}

Retrieval precision is measured using RAGAS \texttt{ContextPrecisionWithoutReference}, which evaluates whether retrieved documents are relevant to the query using automated LLM-based scoring. The RAGAS evaluation judge uses Groq's \texttt{llama-3.3-70b-versatile} model (70B parameters) via the OpenAI-compatible API, providing substantially more capable structured evaluation output than smaller local models.

\subsection{Faithfulness and Hallucination Verification}

Faithfulness is assessed through two complementary mechanisms: (1) RAGAS faithfulness scoring, which decomposes responses into atomic claims and verifies each against retrieved context; and (2) runtime citation grounding ratio, measuring the fraction of response sentences containing valid \texttt{[Source N]} citations.

\subsection{Latency}

End-to-end latency is measured per query, with aggregate statistics computed: mean, median, p95, maximum, and minimum. Latency is decomposed by retrieval mode and query category. The target threshold is p95 $< 2$ seconds for the retrieval component (excluding LLM generation time, which varies by backend).

\subsection{Cost Per Query}

Cost per query is estimated based on prompt and completion token counts, with per-token pricing applied for the active LLM backend. Aggregate monthly cost projections are computed from evaluation workload patterns.

\subsection{Reranking Weight Optimization}

Grid search is performed over the reranking weight space to identify the optimal configuration:

\begin{itemize}[leftmargin=0.5in]
  \item $\alpha \in \{0.4, 0.5, 0.6, 0.7\}$ (semantic similarity weight)
  \item $\beta = \gamma = (1 - \alpha) / 2$ (equal keyword overlap and recency weights)
  \item Constraint: p95 latency $\leq 2.0$ seconds
  \item Objective: maximize mean context precision
\end{itemize}

The configuration yielding the highest mean precision without violating the latency threshold is selected as the production weight configuration.

\begin{table}[H]
\caption{Summary of Evaluation Metrics and Targets.}
\label{tab:eval_targets}
\centering
\begin{tabular}{p{3.0cm}p{3.0cm}p{3.2cm}p{3.0cm}}
\toprule
\textbf{Metric} & \textbf{Dimension} & \textbf{Method} & \textbf{Target} \\
\midrule
Context Precision & Retrieval quality & RAGAS automated & $> 0.75$ \\
Faithfulness & Grounding quality & RAGAS + citation ratio & High (quantitative) \\
p95 Latency & System performance & Per-query timing & $< 2$ seconds \\
Cost per Query & Operational cost & Token usage metrics & $<$~USD 15/month \\
\bottomrule
\end{tabular}
\end{table}

\section{Planned Analysis and Discussion Framework}
\label{sec:discussion}

Upon completion of the full evaluation run, the following analyses will be conducted:

\textbf{Precision vs.\ computational overhead.} Context precision scores will be compared between baseline and hybrid configurations using paired differences across the 25 evaluation queries. The mean precision improvement will be reported alongside the additional per-query latency introduced by the metadata filtering and reranking stages.

\textbf{Retrieval strictness vs.\ recall.} The 180-day metadata filter and minimum similarity threshold (0.15) reduce retrieval noise but may exclude relevant older documents. The analysis will examine whether filtered-out documents would have contributed positively to generation quality by inspecting borderline cases.

\textbf{Cost--performance balance across LLM backends.} Token usage and estimated cost per query will be reported for the Groq backend (primary). Since the prompt template is identical across backends, any faithfulness differences can be attributed to model capability rather than retrieval quality.

\textbf{Reranking weight sensitivity.} Grid search results across $\alpha \in \{0.4, 0.5, 0.6, 0.7\}$ will be presented as a precision--latency Pareto frontier to identify the optimal weight configuration.

\textbf{Production suitability.} The analysis will assess whether all four evaluation targets (context precision $> 0.75$, high faithfulness, p95 latency $< 2$s, cost $<$ USD~15/month) are simultaneously achievable under Learner Lab constraints.

\section{Planned Result Visualizations}
\label{sec:visualizations}

This section presents the planned visualization formats for presenting evaluation results. Placeholder charts illustrate the structure; actual data will be populated upon completion of the full evaluation run (Phase~B).

\subsection{Baseline vs.\ Hybrid Retrieval Comparison}

Figure~\ref{fig:viz_comparison} will present a grouped bar chart comparing baseline and hybrid retrieval across four core evaluation metrics: context precision, faithfulness, citation grounding ratio, and normalized latency. This visualization provides the primary evidence for answering Research Question~2 (Section~\ref{sec:rq}).

\begin{figure}[H]
\centering
\begin{tikzpicture}
\begin{axis}[
    ybar,
    width=0.92\textwidth,
    height=6.5cm,
    bar width=14pt,
    ylabel={Score},
    ylabel style={font=\small},
    xlabel style={font=\small},
    symbolic x coords={Context Precision, Faithfulness, Citation Grounding, Norm.\ Latency},
    xtick=data,
    xticklabel style={font=\scriptsize, rotate=15, anchor=east},
    yticklabel style={font=\scriptsize},
    ymin=0, ymax=1.05,
    legend style={at={(0.5,1.05)}, anchor=south, legend columns=2, font=\small},
    nodes near coords,
    nodes near coords style={font=\tiny, color=black},
    every node near coord/.append style={yshift=2pt},
    grid=major,
    grid style={gray!20},
    area legend,
]
\addplot[fill=awsblue!60, draw=awsblue] coordinates {
    (Context Precision, 0.00)
    (Faithfulness, 0.00)
    (Citation Grounding, 0.00)
    (Norm.\ Latency, 0.00)
};
\addplot[fill=awsorange!60, draw=awsorange] coordinates {
    (Context Precision, 0.00)
    (Faithfulness, 0.00)
    (Citation Grounding, 0.00)
    (Norm.\ Latency, 0.00)
};
\legend{Baseline (Vector-Only), Hybrid (Implemented)}
\end{axis}

\node[font=\small\itshape, color=propgray, anchor=center] at (6.2, 2.5)
  {Data pending --- Phase B evaluation};
\end{tikzpicture}
\caption{Planned comparison of baseline vs.\ hybrid retrieval across four evaluation dimensions. Bars represent mean scores across 25 evaluation queries. Latency is normalized to the 2-second target (value of 1.0 = target met). Actual values will be populated after evaluation execution.}
\label{fig:viz_comparison}
\end{figure}

\subsection{Grid Search: Precision--Latency Pareto Frontier}

Figure~\ref{fig:viz_pareto} will present the grid search results as a scatter plot with the reranking similarity weight ($\alpha$) on the x-axis, context precision on the left y-axis, and p95 latency on the right y-axis. The latency constraint ($\leq 2.0$s) is shown as a horizontal threshold line, enabling identification of the optimal weight configuration that maximizes precision without violating the latency bound.

\begin{figure}[H]
\centering
\begin{tikzpicture}
\begin{axis}[
    width=0.88\textwidth,
    height=6.5cm,
    xlabel={Similarity Weight ($\alpha$)},
    ylabel={Mean Context Precision},
    xlabel style={font=\small},
    ylabel style={font=\small, color=awsblue},
    xmin=0.35, xmax=0.75,
    ymin=0.0, ymax=1.0,
    xtick={0.4, 0.5, 0.6, 0.7},
    xticklabel style={font=\scriptsize},
    yticklabel style={font=\scriptsize, color=awsblue},
    axis y line*=left,
    grid=major,
    grid style={gray!20},
    legend style={at={(0.02,0.98)}, anchor=north west, font=\small},
]
\addplot[color=awsblue, mark=square*, mark size=3pt, line width=1.2pt]
    coordinates {(0.4, 0.0) (0.5, 0.0) (0.6, 0.0) (0.7, 0.0)};
\addlegendentry{Context Precision}
\end{axis}

\begin{axis}[
    width=0.88\textwidth,
    height=6.5cm,
    ylabel={p95 Latency (seconds)},
    ylabel style={font=\small, color=awsorange},
    xmin=0.35, xmax=0.75,
    ymin=0.0, ymax=3.0,
    xtick=\empty,
    yticklabel style={font=\scriptsize, color=awsorange},
    axis y line*=right,
    axis x line=none,
    legend style={at={(0.98,0.98)}, anchor=north east, font=\small},
]
\addplot[color=awsorange, mark=triangle*, mark size=3pt, line width=1.2pt, dashed]
    coordinates {(0.4, 0.0) (0.5, 0.0) (0.6, 0.0) (0.7, 0.0)};
\addlegendentry{p95 Latency}
\draw[propred, thick, dashed] (axis cs:0.35, 2.0) -- (axis cs:0.75, 2.0)
    node[pos=0.85, above, font=\scriptsize, color=propred] {Latency Target ($\leq 2.0$s)};
\end{axis}

\node[font=\small\itshape, color=propgray, anchor=center] at (6.0, 2.8)
  {Data pending --- Phase B grid search};
\end{tikzpicture}
\caption{Planned grid search results: reranking similarity weight ($\alpha$) vs.\ mean context precision (left axis, solid) and p95 latency (right axis, dashed). The horizontal dashed line marks the 2-second latency constraint. The optimal $\alpha$ maximizes precision without exceeding the latency threshold. Constraint: $\beta = \gamma = (1-\alpha)/2$.}
\label{fig:viz_pareto}
\end{figure}

\subsection{Per-Source Retrieval Distribution}

Figure~\ref{fig:viz_sources} will present a stacked bar chart showing the distribution of retrieved documents by source across baseline and hybrid retrieval modes. This visualization reveals whether hybrid retrieval promotes greater source diversity---a key expected benefit of metadata-aware retrieval.

\begin{figure}[H]
\centering
\begin{tikzpicture}
\begin{axis}[
    ybar stacked,
    width=0.7\textwidth,
    height=6.5cm,
    bar width=30pt,
    ylabel={Proportion of Retrieved Documents},
    ylabel style={font=\small},
    symbolic x coords={Baseline, Hybrid},
    xtick=data,
    xticklabel style={font=\small},
    yticklabel style={font=\scriptsize},
    ymin=0, ymax=1.05,
    legend style={at={(1.02,0.5)}, anchor=west, font=\scriptsize, legend columns=1},
    grid=major,
    grid style={gray!20},
    area legend,
]
\addplot[fill=awsblue!70, draw=awsblue] coordinates {(Baseline, 0.20) (Hybrid, 0.20)};
\addplot[fill=awsgreen!60, draw=awsgreen] coordinates {(Baseline, 0.20) (Hybrid, 0.20)};
\addplot[fill=awsorange!60, draw=awsorange] coordinates {(Baseline, 0.20) (Hybrid, 0.20)};
\addplot[fill=awspurple!60, draw=awspurple] coordinates {(Baseline, 0.20) (Hybrid, 0.20)};
\addplot[fill=propgray!50, draw=propgray] coordinates {(Baseline, 0.20) (Hybrid, 0.20)};
\legend{ArXiv, Hacker News, DEV.to, GitHub, RSS}
\end{axis}

\node[font=\small\itshape, color=propgray, anchor=center] at (4.0, 2.8)
  {Uniform placeholder --- actual distribution pending};
\end{tikzpicture}
\caption{Planned per-source retrieval distribution for baseline vs.\ hybrid modes. Stacked bars show the proportion of top-$k$ retrieved documents originating from each of the five data sources, averaged across 25 evaluation queries. Greater diversity in the hybrid configuration would support the retrieval hypothesis.}
\label{fig:viz_sources}
\end{figure}

\subsection{Latency Breakdown by Pipeline Stage}

Figure~\ref{fig:viz_latency} will present a horizontal stacked bar chart decomposing end-to-end latency into four pipeline stages: metadata filtering, vector search, reranking, and LLM generation. This enables identification of the primary latency bottleneck and validates the ``lightweight'' characterization of the reranking stage.

\begin{figure}[H]
\centering
\begin{tikzpicture}
\begin{axis}[
    xbar stacked,
    width=0.85\textwidth,
    height=5.0cm,
    bar width=22pt,
    xlabel={Mean Latency (seconds)},
    xlabel style={font=\small},
    ytick={0,1},
    yticklabels={Baseline,Hybrid},
    yticklabel style={font=\small},
    xticklabel style={font=\scriptsize},
    xmin=0, xmax=3.0,
    ymin=-0.6, ymax=1.6,
    legend style={at={(0.5,-0.3)}, anchor=north, legend columns=4, font=\scriptsize},
    grid=major,
    grid style={gray!20},
    area legend,
]
\addplot[fill=propblue!50, draw=propblue] coordinates {(0.0,0) (0.0,1)};
\addplot[fill=awsblue!50, draw=awsblue] coordinates {(0.0,0) (0.0,1)};
\addplot[fill=awspurple!50, draw=awspurple] coordinates {(0.0,0) (0.0,1)};
\addplot[fill=awsorange!50, draw=awsorange] coordinates {(0.0,0) (0.0,1)};
\legend{Metadata Filter, Vector Search, Reranking, LLM Generation}

\draw[propred, thick, dashed] (axis cs:2.0,-0.6) -- (axis cs:2.0,1.6)
    node[pos=0.85, right, font=\scriptsize, color=propred, xshift=2pt] {p95 Target};
\end{axis}

\node[font=\small\itshape, color=propgray, anchor=center] at (5.5, 2.0)
  {Data pending --- Phase B profiling};
\end{tikzpicture}
\caption{Planned latency decomposition by pipeline stage for baseline and hybrid retrieval. The baseline omits metadata filtering and reranking stages. The vertical dashed line marks the 2-second p95 latency target. Actual values will be populated from per-query timing measurements.}
\label{fig:viz_latency}
\end{figure}

\subsection{Evaluation Summary Dashboard}

Table~\ref{tab:viz_dashboard} will serve as the primary summary dashboard, consolidating all evaluation metrics into a single comparison table with target thresholds and pass/fail status indicators.

\begin{table}[H]
\caption{Planned Evaluation Summary Dashboard: Baseline vs.\ Hybrid with Target Thresholds.}
\label{tab:viz_dashboard}
\centering
\begin{tabular}{p{3.5cm}p{2.0cm}p{2.0cm}p{2.0cm}p{2.5cm}}
\toprule
\textbf{Metric} & \textbf{Target} & \textbf{Baseline} & \textbf{Hybrid} & \textbf{Status} \\
\midrule
Context Precision & $> 0.75$ & --- & --- & \textit{pending} \\
Faithfulness & High & --- & --- & \textit{pending} \\
Citation Grounding & $> 0.80$ & --- & --- & \textit{pending} \\
p95 Latency & $< 2.0$s & --- & --- & \textit{pending} \\
Mean Latency & --- & --- & --- & \textit{pending} \\
Cost / Month & $<$ USD 15 & --- & --- & \textit{pending} \\
Optimal $\alpha$ & --- & N/A & --- & \textit{pending} \\
Source Diversity (entropy) & --- & --- & --- & \textit{pending} \\
\bottomrule
\end{tabular}
\end{table}

% ============================================================
% CHAPTER 6: IMPLEMENTATION PROGRESS
% ============================================================
\chapter{Implementation Progress and Deviations}
\label{ch:progress}

This chapter documents the current implementation status, deviations from the original proposal, and key engineering decisions made during development.

\section{Overall Implementation Status}
\label{sec:impl_status}

Table~\ref{tab:impl_status} summarizes the implementation status of each system component as of the reporting date.

\begin{table}[H]
\caption{Component Implementation Status Summary.}
\label{tab:impl_status}
\centering
\begin{tabular}{p{4.5cm}p{2.0cm}p{6.0cm}}
\toprule
\textbf{Component} & \textbf{Status} & \textbf{Notes} \\
\midrule
Data Ingestion (5 sources) & Complete & ArXiv, HN, DEV.to, GitHub, RSS fully implemented \\
Preprocessing Pipeline & Complete & 17-step normalization + token-based chunking \\
Embedding Pipeline & Complete & MiniLM L6 v2 (384-dim), lazy singleton \\
Hybrid Retrieval & Complete & 3-stage pipeline with configurable weights \\
Baseline Retrieval & Complete & Vector-only comparison mode \\
RAG Orchestrator & Complete & 3-layer hallucination verification \\
LLM Backend Fallback & Complete & Groq (primary) $\to$ Bedrock $\to$ Ollama $\to$ HuggingFace \\
API Layer & Complete & FastAPI + Mangum; /health, /ask, /drift endpoints \\
Database Schema & Complete & PostgreSQL + pgvector with HNSW index \\
S3 Medallion Storage & Complete & Raw/processed/embeddings tiers \\
SQS Queue + DLQ & Complete & Ingestion--embedding decoupling \\
Drift Detection & Complete & Dual-criteria: 10\% threshold + Shewhart $3\sigma$ \\
Health Monitoring & Complete & 4-probe deep health check (DB, S3, SQS, LLM) \\
CloudWatch Integration & Complete & 6 custom metric namespaces \\
Evaluation Framework & Complete & 25 queries, RAGAS, grid search (code complete; full run pending) \\
CI/CD Pipeline & Complete & 4-stage GitHub Actions + SAM \\
SAM Infrastructure (Learner Lab) & Complete & Full SAM template with 15+ AWS resources \\
Docker Compose (Local) & Complete & 6 services (db, pgadmin, api, frontend, scheduler, localstack) \\
React Frontend & Complete & Query interface with citation display \\
Unit Tests & Complete & 49 tests across 5 test files, all passing \\
\bottomrule
\end{tabular}
\end{table}

\section{Deviations from the Original Proposal}
\label{sec:deviations}

Several implementation decisions diverge from the original proposal due to AWS Academy Learner Lab constraints and practical engineering considerations. Table~\ref{tab:deviations} documents each deviation, its rationale, and the implemented alternative.

\begin{table}[H]
\caption{Deviations from the Original Proposal and Rationale.}
\label{tab:deviations}
\centering
\begin{tabular}{p{3.0cm}p{3.0cm}p{3.0cm}p{3.5cm}}
\toprule
\textbf{Proposal Claim} & \textbf{Implemented} & \textbf{Reason} & \textbf{Impact} \\
\midrule
Aurora Serverless + pgvector & RDS PostgreSQL db.t3.micro + pgvector & Aurora not available in Learner Lab & Functionally equivalent; HNSW index preserved \\[4pt]
Amazon Bedrock (Claude Haiku) & Groq (llama-3.1-8b-instant) as primary + auto-fallback chain & Educational AWS account has 0 Bedrock model quota (not adjustable); Groq free tier provides reliable generation & Same prompt template; evaluation comparable; Groq latency $<$ Bedrock \\[4pt]
2 data sources (ArXiv, HN) & 5 data sources (+DEV.to, GitHub, RSS) & Improved corpus diversity and breadth & Enhanced retrieval coverage \\[4pt]
Custom VPC + NAT Gateway & Default VPC & Custom VPCs restricted in Learner Lab & Simplified networking; publicly accessible RDS \\[4pt]
AWS Budgets alerts & Cost Explorer programmatic check & AWS Budgets restricted in Learner Lab & Budget guard preserved via Cost Explorer API \\[4pt]
CDK or SAM deployment & SAM only & Simplified IaC management & Reduced toolchain complexity \\[4pt]
HN every 5 minutes & HN every 30 minutes & Rate limit considerations & Reduced API pressure \\[4pt]
ArXiv weekly & ArXiv every 12 hours & More frequent scholarly updates & Improved corpus freshness \\[4pt]
Simple 10\% drift threshold & Dual: 10\% threshold + Shewhart $3\sigma$ & Stronger statistical process control & More robust drift detection \\[4pt]
Single LLM backend selection & Automatic fallback chain & Improve availability & System resilience \\
\bottomrule
\end{tabular}
\end{table}

\section{Expanded Data Source Coverage}
\label{sec:expanded_sources}

The most significant enhancement beyond the original proposal is the expansion from two to five data sources. This decision was motivated by the need for broader topical coverage and diverse temporal profiles. The three additional sources---DEV.to (practitioner tutorials across 21 technology tags), GitHub Trending (emerging open-source projects across 10 topic queries), and curated RSS/Atom feeds (10 technology journalism outlets)---complement the scholarly authority of ArXiv and the community curation of Hacker News.

All five ingesters share the same idempotent, SHA-256-based deduplication architecture, ensuring that the expanded ingestion scope does not compromise data integrity.

\section{Enhanced Preprocessing Pipeline}
\label{sec:enhanced_preprocessing}

The preprocessing pipeline was substantially enhanced beyond the original proposal specification. The initial design specified HTML stripping, encoding normalization, whitespace consolidation, and removal of low-information tokens. The implemented 17-step normalization pipeline adds handling for: emoji and pictographic symbols (6 Unicode ranges), Unicode control characters and zero-width markers, Markdown blockquotes, HTML comments, strikethrough syntax, and repeated punctuation sequences. These additions were motivated by the diversity of content formats encountered across the five data sources, particularly the Markdown-heavy content from DEV.to and GitHub README files. The preprocessing pipeline is validated by 18 dedicated unit tests covering each normalization step.

\section{Learner Lab Adaptation}
\label{sec:learnerlab}

Adapting the system for AWS Academy Learner Lab required systematic identification and replacement of restricted services while preserving the core architectural capabilities. The key constraints and adaptations are:

\begin{itemize}[leftmargin=0.5in]
  \item \textbf{No Aurora Serverless.} Replaced with standard RDS PostgreSQL \texttt{db.t3.micro} (PostgreSQL 16.4, 20GB gp2). The pgvector extension and HNSW indexing remain fully functional.
  \item \textbf{Groq as primary LLM backend.} The educational AWS account (AWS Academy) has all Amazon Bedrock model quotas set to 0 and marked ``Not adjustable,'' making Bedrock unusable for generation. The system now uses Groq (\texttt{LLM\_BACKEND=groq}) as the primary LLM backend, running \texttt{llama-3.1-8b-instant} via an OpenAI-compatible API with a free tier sufficient for the project workload. Bedrock, Ollama, and HuggingFace remain in the automatic fallback chain. The RAGAS evaluation judge uses a larger Groq model (\texttt{llama-3.3-70b-versatile}, 70B parameters) for more reliable structured scoring output.
  \item \textbf{No custom VPCs or NAT Gateways.} The SAM template uses the default VPC with two subnets. RDS is configured as publicly accessible with security group restrictions.
  \item \textbf{No custom IAM roles.} All Lambda functions use the pre-existing \texttt{LabRole} provided by the Learner Lab environment.
  \item \textbf{No AWS Budgets service.} Budget governance is maintained through programmatic Cost Explorer API queries integrated into the inference path.
  \item \textbf{LLM backend set to \texttt{groq}.} The \texttt{LLM\_BACKEND} environment variable is set to \texttt{groq} in the Free Tier SAM template, with \texttt{BUDGET\_HALT\_ENABLED} set to \texttt{false} (Cost Explorer quota limitations). The Groq API key is stored as a GitHub Actions secret (\texttt{GROQ\_API\_KEY}) and injected at deploy time.
\end{itemize}

A dedicated SAM template (\texttt{template-learnerlab.yaml}) and configuration file (\texttt{samconfig-learnerlab.toml}) were created to encapsulate all Learner Lab-specific resource definitions.

\section{Test Coverage}
\label{sec:tests}

The system includes 49 unit tests distributed across five test modules:

\begin{table}[H]
\caption{Unit Test Distribution.}
\label{tab:tests}
\centering
\begin{tabular}{p{4.5cm}p{1.5cm}p{6.5cm}}
\toprule
\textbf{Test Module} & \textbf{Count} & \textbf{Coverage} \\
\midrule
\texttt{test\_preprocessing.py} & 18 & All 17 normalization steps + chunking behavior \\
\texttt{test\_ingestion.py} & 14 & SHA-256 hash determinism and uniqueness for all 5 sources \\
\texttt{test\_retriever.py} & 8 & Keyword overlap computation + recency weight decay \\
\texttt{test\_rag.py} & 7 & Citation grounding check + context block construction \\
\texttt{test\_api.py} & 5 & FastAPI endpoint validation (health, error handling, success) \\
\midrule
\textbf{Total} & \textbf{49} (passing) & \\
\bottomrule
\end{tabular}
\end{table}

All tests execute without external dependencies (database, AWS services) using mocked interfaces where appropriate.

\section{Custom CloudWatch Metrics}
\label{sec:metrics}

Six custom CloudWatch metric namespaces have been implemented to provide operational visibility:

\begin{enumerate}[leftmargin=0.5in]
  \item \texttt{IngestionCount} --- documents ingested per source per cycle
  \item \texttt{PipelineChunks} --- chunks generated per pipeline run
  \item \texttt{PipelineLatency} --- end-to-end pipeline execution time
  \item \texttt{PipelineErrors} --- pipeline failure count
  \item \texttt{ApiLatency} --- per-request API response time
  \item \texttt{HallucinationFlags} --- count of citation grounding failures
\end{enumerate}

Additionally, the drift detection module publishes \texttt{DriftMeanSimilarity} and \texttt{DriftAlert} metrics for time-series tracking of retrieval quality.

\section{System Demonstration}
\label{sec:demo}

This section presents screenshots from the running system deployed locally via Docker Compose, demonstrating end-to-end functionality across all implemented components.

\subsection{Frontend: Hybrid Retrieval Query with Citation-Grounded Response}

Figure~\ref{fig:demo_hybrid} shows the TechPulse frontend answering a technology trend query using the hybrid retrieval mode. The response includes inline \texttt{[Source N]} citations and a source list with relevance scores, demonstrating the complete RAG pipeline from query submission through grounded generation.

\begin{figure}[H]
\centering
\fbox{\parbox{0.9\textwidth}{\centering\vspace{3cm}
\textit{Screenshot: TechPulse frontend --- hybrid mode query with citation-grounded answer and source list.}
\\[4pt]
\textbf{File to place:} \texttt{figures/demo\_hybrid.png}
\vspace{3cm}}}
\caption{TechPulse frontend: hybrid retrieval query with citation-grounded response. The response displays inline source citations and a ranked source list with relevance scores from all five data sources.}
\label{fig:demo_hybrid}
\end{figure}

\subsection{Frontend: Baseline vs.\ Hybrid Comparison}

Figure~\ref{fig:demo_baseline} shows the same query executed in baseline (vector-only) mode for comparison. The difference in source diversity and citation quality between baseline and hybrid modes provides qualitative evidence supporting the hybrid retrieval hypothesis.

\begin{figure}[H]
\centering
\fbox{\parbox{0.9\textwidth}{\centering\vspace{3cm}
\textit{Screenshot: TechPulse frontend --- baseline mode query on the same question for comparison.}
\\[4pt]
\textbf{File to place:} \texttt{figures/demo\_baseline.png}
\vspace{3cm}}}
\caption{TechPulse frontend: baseline (vector-only) retrieval query on the same question as Figure~\ref{fig:demo_hybrid}, enabling visual comparison of response quality and source diversity.}
\label{fig:demo_baseline}
\end{figure}

\subsection{Unit Test Execution}

Figure~\ref{fig:demo_tests} shows the terminal output of running the full test suite. All 49 tests pass across the five test modules, validating preprocessing, ingestion, retrieval, RAG orchestration, and API endpoint behavior.

\begin{figure}[H]
\centering
\fbox{\parbox{0.9\textwidth}{\centering\vspace{3cm}
\textit{Screenshot: Terminal output of \texttt{pytest tests/ -v} showing 49/49 tests passing.}
\\[4pt]
\textbf{File to place:} \texttt{figures/demo\_tests.png}
\vspace{3cm}}}
\caption{Pytest execution: 49 unit tests passing across 5 test modules (\texttt{test\_preprocessing}, \texttt{test\_ingestion}, \texttt{test\_retriever}, \texttt{test\_rag}, \texttt{test\_api}).}
\label{fig:demo_tests}
\end{figure}

\subsection{Docker Compose: Running Services}

Figure~\ref{fig:demo_docker} shows all six Docker Compose services running simultaneously: PostgreSQL with pgvector, pgAdmin, LocalStack (S3/SQS), FastAPI backend, React frontend, and the ingestion scheduler.

\begin{figure}[H]
\centering
\fbox{\parbox{0.9\textwidth}{\centering\vspace{3cm}
\textit{Screenshot: \texttt{docker ps} output showing all 6 containers running (db, pgadmin, localstack, api, frontend, scheduler).}
\\[4pt]
\textbf{File to place:} \texttt{figures/demo\_docker.png}
\vspace{3cm}}}
\caption{Docker Compose deployment: six services running locally --- PostgreSQL+pgvector, pgAdmin, LocalStack, FastAPI API, React frontend, and the ingestion scheduler.}
\label{fig:demo_docker}
\end{figure}

\subsection{CI/CD Pipeline: GitHub Actions}

Figure~\ref{fig:demo_cicd} shows the GitHub Actions CI/CD pipeline executing the four-stage workflow: lint-and-test, Docker build, SAM build, and deployment.

\begin{figure}[H]
\centering
\fbox{\parbox{0.9\textwidth}{\centering\vspace{3cm}
\textit{Screenshot: GitHub Actions workflow showing all 4 stages (lint-and-test, docker-build, sam-build, deploy-dev) passing.}
\\[4pt]
\textbf{File to place:} \texttt{figures/demo\_cicd.png}
\vspace{3cm}}}
\caption{GitHub Actions CI/CD pipeline: four-stage workflow (lint \& test, Docker build, SAM build, deploy) executing successfully on the \texttt{main} branch.}
\label{fig:demo_cicd}
\end{figure}

\subsection{Database: Populated Corpus from Five Sources}

Figure~\ref{fig:demo_pgadmin} shows the pgAdmin interface with the \texttt{documents} table containing ingested records from all five data sources, demonstrating that the ingestion pipeline successfully populates the corpus.

\begin{figure}[H]
\centering
\fbox{\parbox{0.9\textwidth}{\centering\vspace{3cm}
\textit{Screenshot: pgAdmin showing the documents table with rows from ArXiv, HN, DEV.to, GitHub, and RSS sources.}
\\[4pt]
\textbf{File to place:} \texttt{figures/demo\_pgadmin.png}
\vspace{3cm}}}
\caption{pgAdmin: the \texttt{documents} table populated with records from all five data sources (ArXiv, Hacker News, DEV.to, GitHub Trending, RSS), demonstrating successful multi-source ingestion.}
\label{fig:demo_pgadmin}
\end{figure}

% ============================================================
% CHAPTER 7: REMAINING WORK AND LIMITATIONS
% ============================================================
\chapter{Remaining Work and Limitations}

\section{Remaining Work}
\label{sec:remaining}

The remaining work is organized into four phases spanning from mid-March through the second week of April 2026. Table~\ref{tab:remaining_phases} provides a phase-level overview, followed by detailed task descriptions.

\begin{table}[H]
\caption{Remaining Work: Phase Overview.}
\label{tab:remaining_phases}
\centering
\begin{tabular}{p{0.8cm}p{4.2cm}p{2.5cm}p{5.0cm}}
\toprule
\textbf{Phase} & \textbf{Focus} & \textbf{Timeline} & \textbf{Key Deliverables} \\
\midrule
A & System Validation \& Corpus Building & Mar W3--W4 & Populated corpus; local integration verified \\
B & Evaluation \& Experimentation & Mar W4--Apr W1 & Empirical results; grid search output \\
C & Analysis \& Monitoring Validation & Apr W1--W2 & Drift validation; results visualization \\
D & Final Report \& Submission & Apr W2 & Completed report with empirical evidence \\
\bottomrule
\end{tabular}
\end{table}

\subsection{Phase A: System Validation and Corpus Building (Mar W3 -- Mar W4)}

\begin{enumerate}[leftmargin=0.5in]
  \item \textbf{Corpus population via multi-cycle ingestion.} Execute scheduled ingestion cycles across all five data sources over a sustained period to build a sufficiently large and diverse corpus. Target: 1,000+ indexed documents with balanced source representation across ArXiv, Hacker News, DEV.to, GitHub Trending, and RSS feeds.

  \item \textbf{Local Docker Compose integration testing.} Verify end-to-end pipeline functionality in the local Docker Compose environment: ingestion $\to$ preprocessing $\to$ embedding $\to$ indexing $\to$ retrieval $\to$ generation. Confirm all six services (db, pgadmin, localstack, api, frontend, scheduler) operate correctly together.

  \item \textbf{Frontend integration verification.} Test the React frontend against the live API to verify citation-grounded response rendering, baseline/hybrid mode toggling, abstention message handling, source badge display, and retrieval metadata transparency.

  \item \textbf{AWS Learner Lab deployment (in progress).} Deploy the full system to the Learner Lab environment using the SAM template (\texttt{template-learnerlab.yaml}). Deployment began in parallel with Docker integration testing; the SAM template and \texttt{samconfig-learnerlab.toml} have been prepared and validated, and initial stack creation is underway. Remaining: verify Lambda function execution, RDS connectivity, S3 storage operations, SQS message flow, and API Gateway routing under cloud conditions. This task runs continuously through the end of the project to support iterative testing and evaluation in the cloud environment.
\end{enumerate}

\subsection{Phase B: Evaluation and Experimentation (Mar W4 -- Apr W1)}

\begin{enumerate}[leftmargin=0.5in, resume]
  \item \textbf{Full baseline evaluation run.} Execute 25 evaluation queries in vector-only (baseline) mode against the populated corpus. Record context precision, faithfulness, citation grounding ratio, and per-query latency for each query.

  \item \textbf{Full hybrid evaluation run.} Execute the same 25 evaluation queries in hybrid mode using the default reranking weights ($\alpha = 0.6$, $\beta = 0.2$, $\gamma = 0.2$). Record all metrics under identical corpus and embedding conditions to enable paired comparison.

  \item \textbf{Grid search weight optimization.} Run the reranking weight grid search across $\alpha \in \{0.4, 0.5, 0.6, 0.7\}$ with the latency constraint (p95 $\leq 2.0$s). Identify the optimal weight configuration that maximizes mean context precision.

  \item \textbf{Latency profiling and cost estimation.} Decompose end-to-end latency by stage (metadata filtering, vector search, reranking, LLM generation). Compute per-query cost estimates based on measured token usage across LLM backends.

  \item \textbf{Source diversity analysis.} Investigate the per-source distribution of retrieved results to determine whether the corpus is skewed toward any single source (e.g., DEV.to). If imbalance is confirmed, evaluate mitigation strategies: balanced ingestion quotas per source, source-aware reranking weight adjustments, or post-retrieval diversity enforcement.

  \item \textbf{HTTPS frontend via CloudFront (optional).} Add a CloudFront distribution in front of the S3 static website to serve the React frontend over HTTPS. This eliminates the browser ``Not secure'' warning visible in the current HTTP deployment and improves the production-readiness appearance for demonstration.
\end{enumerate}

\subsection{Phase C: Analysis and Monitoring Validation (Apr W1 -- Apr W2)}

\begin{enumerate}[leftmargin=0.5in, resume]
  \item \textbf{Drift detection validation.} Execute multiple drift detection cycles against the populated corpus to validate the dual-criteria mechanism (10\% relative threshold + Shewhart $3\sigma$). Confirm that alerts trigger correctly under simulated drift conditions and that baseline history accumulates as expected.

  \item \textbf{Results analysis and visualization.} Produce comparison tables and visualizations: baseline vs.\ hybrid precision paired differences, grid search Pareto frontier (precision vs.\ latency), per-source retrieval distribution, and latency breakdown charts.

  \item \textbf{System demonstration screenshots.} Capture screenshots for all placeholder figures in the report (Figures~\ref{fig:demo_hybrid}--\ref{fig:demo_pgadmin} and Appendix~\ref{app:screenshots}), replacing placeholder boxes with actual system output from both Docker Compose and Learner Lab environments.
\end{enumerate}

\subsection{Phase D: Final Report and Submission (Apr W2)}

\begin{enumerate}[leftmargin=0.5in, resume]
  \item \textbf{Evaluation chapter completion.} Write the complete evaluation results chapter with empirical findings, statistical comparison of baseline vs.\ hybrid retrieval, discussion of grid search outcomes, and analysis against the four evaluation targets (precision $> 0.75$, high faithfulness, p95 $< 2$s, cost $<$ USD~15/month).

  \item \textbf{Discussion and conclusions.} Complete the discussion of results including the five planned analyses (Section~\ref{sec:discussion}): precision vs.\ overhead, retrieval strictness vs.\ recall, cost--performance balance, weight sensitivity, and production suitability assessment.

  \item \textbf{Final report revision and submission.} Integrate all empirical content, screenshots, and analysis into the final report. Conduct internal peer review, address formatting and consistency, and prepare the submission package including source code archive.
\end{enumerate}

\section{Known Limitations}
\label{sec:limitations}

\begin{enumerate}[leftmargin=0.5in]
  \item The evaluation corpus is limited to five technology-focused data sources; generalizability to other domains or larger corpora requires further study.

  \item The evaluation query set of 25 curated queries is sufficient for preliminary comparative analysis consistent with prior RAG evaluation practices (Es et al., 2023). However, a larger and more diverse query set would strengthen external validity.

  \item Amazon Bedrock model quotas are set to 0 on the educational AWS account and marked ``Not adjustable,'' requiring the use of Groq (free tier) as the primary LLM backend. Since the prompt template, evaluation methodology, and RAGAS metrics are backend-agnostic, the comparative findings between baseline and hybrid retrieval remain valid and transferable to any LLM backend. The RAGAS evaluation judge uses Groq's \texttt{llama-3.3-70b-versatile} (70B parameters) to ensure reliable structured evaluation output.

  \item The retrieval results may exhibit source diversity skew---preliminary observations suggest DEV.to articles dominate the retrieved sources due to higher ingestion volume and embedding density relative to other sources. This corpus imbalance may affect evaluation fairness across source types and requires investigation (balanced ingestion quotas or source-aware reranking).

  \item The frontend is served over HTTP via S3 static website hosting. Adding a CloudFront distribution would provide HTTPS encryption and improve the production-readiness appearance for demonstration and submission.

  \item The \texttt{db.t3.micro} RDS instance provides limited compute and memory compared to Aurora Serverless. Under high ingestion volumes, database performance may become a bottleneck.

  \item Cross-encoder reranking and ANN-optimized indexing strategies are out of scope for this project but identified as future research directions.

  \item Budget governance via Cost Explorer has higher latency than the originally proposed AWS Budgets integration, as Cost Explorer data may lag by up to 24 hours.

  \item The current unit test suite (49 tests) covers core logic but lacks integration tests that exercise the full pipeline end-to-end with a live database and AWS services.
\end{enumerate}

\section{Project Timeline}
\label{sec:timeline}

Figure~\ref{fig:gantt} presents the project timeline, including completed milestones and remaining work.

\begin{figure}[H]
\centering
\resizebox{\textwidth}{!}{%
\begin{tikzpicture}

% Header — 8 time columns (Apr W3/W4 removed; deadline = Apr W2)
\node[font=\small\bfseries] at (-4.0, 0.5) {Phase};
\foreach \m/\lab [count=\i from 0] in {
  1/{Jan},2/{Feb},3/{Mar-W1},4/{Mar-W2},5/{Mar-W3},6/{Mar-W4},7/{Apr-W1},8/{Apr-W2}} {
  \node[font=\scriptsize\bfseries] at ({\i*1.5+0.75}, 0.5) {\lab};
  \draw[awsgray!20] ({\i*1.5}, -15.2) -- ({\i*1.5}, 0.2);
}

\def\barh{0.40}
\def\rsp{-1.0}  % row spacing

% ================================================================
% COMPLETED ROWS (solid fill)
% ================================================================

% Row 1
\node[font=\scriptsize, anchor=east] at (-0.2, {1*\rsp+0.5}) {Requirements \& Design};
\fill[awsblue!60] (0,{1*\rsp+0.5-\barh/2}) rectangle (3.0,{1*\rsp+0.5+\barh/2});
\node[font=\scriptsize, white] at (1.5, {1*\rsp+0.5}) {\checkmark};

% Row 2
\node[font=\scriptsize, anchor=east] at (-0.2, {2*\rsp+0.5}) {Data Ingestion (5 sources)};
\fill[awsgreen!60] (1.5,{2*\rsp+0.5-\barh/2}) rectangle (6.0,{2*\rsp+0.5+\barh/2});
\node[font=\scriptsize, white] at (3.75, {2*\rsp+0.5}) {\checkmark};

% Row 3
\node[font=\scriptsize, anchor=east] at (-0.2, {3*\rsp+0.5}) {Preprocessing \& Embedding};
\fill[awsgreen!60] (3.0,{3*\rsp+0.5-\barh/2}) rectangle (6.0,{3*\rsp+0.5+\barh/2});
\node[font=\scriptsize, white] at (4.5, {3*\rsp+0.5}) {\checkmark};

% Row 4
\node[font=\scriptsize, anchor=east] at (-0.2, {4*\rsp+0.5}) {Hybrid Retrieval \& RAG};
\fill[awspurple!60] (3.0,{4*\rsp+0.5-\barh/2}) rectangle (7.5,{4*\rsp+0.5+\barh/2});
\node[font=\scriptsize, white] at (5.25, {4*\rsp+0.5}) {\checkmark};

% Row 5
\node[font=\scriptsize, anchor=east] at (-0.2, {5*\rsp+0.5}) {Infrastructure (SAM/CI-CD)};
\fill[awsorange!60] (4.5,{5*\rsp+0.5-\barh/2}) rectangle (7.5,{5*\rsp+0.5+\barh/2});
\node[font=\scriptsize, white] at (6.0, {5*\rsp+0.5}) {\checkmark};

% Row 6
\node[font=\scriptsize, anchor=east] at (-0.2, {6*\rsp+0.5}) {Evaluation Framework (code)};
\fill[awsgreen!60] (4.5,{6*\rsp+0.5-\barh/2}) rectangle (7.5,{6*\rsp+0.5+\barh/2});
\node[font=\scriptsize, white] at (6.0, {6*\rsp+0.5}) {\checkmark};

% ================================================================
% PHASE A — System Validation & Corpus Building (Mar W3 – Mar W4)
% ================================================================
\node[font=\scriptsize\itshape, color=awsgray, anchor=east] at (-0.2, {6.7*\rsp+0.5})
  {\textbf{Phase A}};

% Row 7
\node[font=\scriptsize, anchor=east] at (-0.2, {7*\rsp+0.5}) {Corpus Population};
\fill[propblue!25] (6.0,{7*\rsp+0.5-\barh/2}) rectangle (9.0,{7*\rsp+0.5+\barh/2});
\draw[propblue!60, thick] (6.0,{7*\rsp+0.5-\barh/2}) rectangle (9.0,{7*\rsp+0.5+\barh/2});

% Row 8
\node[font=\scriptsize, anchor=east] at (-0.2, {8*\rsp+0.5}) {Docker Integration Testing};
\fill[propblue!25] (6.0,{8*\rsp+0.5-\barh/2}) rectangle (7.5,{8*\rsp+0.5+\barh/2});
\draw[propblue!60, thick] (6.0,{8*\rsp+0.5-\barh/2}) rectangle (7.5,{8*\rsp+0.5+\barh/2});

% Row 9  (starts same time as Docker testing, extends to end)
\node[font=\scriptsize, anchor=east] at (-0.2, {9*\rsp+0.5}) {AWS Learner Lab Deploy};
\fill[propblue!25] (6.0,{9*\rsp+0.5-\barh/2}) rectangle (12.0,{9*\rsp+0.5+\barh/2});
\draw[propblue!60, thick] (6.0,{9*\rsp+0.5-\barh/2}) rectangle (12.0,{9*\rsp+0.5+\barh/2});

% ================================================================
% PHASE B — Evaluation & Experimentation (Mar W4 – Apr W1)
% ================================================================
\node[font=\scriptsize\itshape, color=awsgray, anchor=east] at (-0.2, {9.7*\rsp+0.5})
  {\textbf{Phase B}};

% Row 10
\node[font=\scriptsize, anchor=east] at (-0.2, {10*\rsp+0.5}) {Evaluation Run (baseline + hybrid)};
\fill[awsorange!20] (7.5,{10*\rsp+0.5-\barh/2}) rectangle (10.5,{10*\rsp+0.5+\barh/2});
\draw[awsorange!60, thick] (7.5,{10*\rsp+0.5-\barh/2}) rectangle (10.5,{10*\rsp+0.5+\barh/2});

% Row 11
\node[font=\scriptsize, anchor=east] at (-0.2, {11*\rsp+0.5}) {Grid Search \& Latency Profiling};
\fill[awsorange!20] (9.0,{11*\rsp+0.5-\barh/2}) rectangle (10.5,{11*\rsp+0.5+\barh/2});
\draw[awsorange!60, thick] (9.0,{11*\rsp+0.5-\barh/2}) rectangle (10.5,{11*\rsp+0.5+\barh/2});

% ================================================================
% PHASE C — Analysis & Monitoring Validation (Apr W1 – Apr W2)
% ================================================================
\node[font=\scriptsize\itshape, color=awsgray, anchor=east] at (-0.2, {11.7*\rsp+0.5})
  {\textbf{Phase C}};

% Row 12
\node[font=\scriptsize, anchor=east] at (-0.2, {12*\rsp+0.5}) {Drift Validation \& Monitoring};
\fill[awspurple!20] (9.0,{12*\rsp+0.5-\barh/2}) rectangle (10.5,{12*\rsp+0.5+\barh/2});
\draw[awspurple!60, thick] (9.0,{12*\rsp+0.5-\barh/2}) rectangle (10.5,{12*\rsp+0.5+\barh/2});

% Row 13
\node[font=\scriptsize, anchor=east] at (-0.2, {13*\rsp+0.5}) {Results Analysis \& Visualization};
\fill[awspurple!20] (9.0,{13*\rsp+0.5-\barh/2}) rectangle (12.0,{13*\rsp+0.5+\barh/2});
\draw[awspurple!60, thick] (9.0,{13*\rsp+0.5-\barh/2}) rectangle (12.0,{13*\rsp+0.5+\barh/2});

% ================================================================
% PHASE D — Final Report & Submission (Apr W2)
% ================================================================
\node[font=\scriptsize\itshape, color=awsgray, anchor=east] at (-0.2, {13.7*\rsp+0.5})
  {\textbf{Phase D}};

% Row 14
\node[font=\scriptsize, anchor=east] at (-0.2, {14*\rsp+0.5}) {Final Report \& Submission};
\fill[propred!25] (10.5,{14*\rsp+0.5-\barh/2}) rectangle (12.0,{14*\rsp+0.5+\barh/2});
\draw[propred!60, thick] (10.5,{14*\rsp+0.5-\barh/2}) rectangle (12.0,{14*\rsp+0.5+\barh/2});

% ================================================================
% Current date marker — Mar W3 midpoint ≈ x=6.75
% ================================================================
\draw[propred, very thick, dashed] (6.75, 0.8) -- (6.75, -14.5);
\node[font=\scriptsize\bfseries, propred] at (6.75, 1.1) {Current Status};

% ================================================================
% Legend
% ================================================================
\fill[awsblue!60] (0.5,-17.2) rectangle (1.3,-16.9);
\node[font=\scriptsize, anchor=west] at (1.4, -17.05) {Completed};

\fill[propblue!25] (3.5,-17.2) rectangle (4.3,-16.9);
\draw[propblue!60, thick] (3.5,-17.2) rectangle (4.3,-16.9);
\node[font=\scriptsize, anchor=west] at (4.4, -17.05) {Phase A (Validation)};

\fill[awsorange!20] (7.0,-17.2) rectangle (7.8,-16.9);
\draw[awsorange!60, thick] (7.0,-17.2) rectangle (7.8,-16.9);
\node[font=\scriptsize, anchor=west] at (7.9, -17.05) {Phase B (Evaluation)};

\fill[awspurple!20] (0.5,-17.7) rectangle (1.3,-17.4);
\draw[awspurple!60, thick] (0.5,-17.7) rectangle (1.3,-17.4);
\node[font=\scriptsize, anchor=west] at (1.4, -17.55) {Phase C};

\fill[propred!25] (3.5,-17.7) rectangle (4.3,-17.4);
\draw[propred!60, thick] (3.5,-17.7) rectangle (4.3,-17.4);
\node[font=\scriptsize, anchor=west] at (4.4, -17.55) {Phase D (Final Report)};

\end{tikzpicture}%
}
\caption{Project timeline (Gantt chart). Solid bars indicate completed phases; outlined bars indicate remaining work organized into four phases (A--D). The dashed vertical line marks the progress report submission. All remaining work is scheduled to complete by second week of April 2026.}
\label{fig:gantt}
\end{figure}

\section{Team Contributions}
\label{sec:contributions}

Table~\ref{tab:contributions} summarises the primary responsibilities of each team member.

\begin{table}[H]
\caption{Team Member Contribution Summary.}
\label{tab:contributions}
\centering
\begin{tabular}{p{4.0cm}p{4.0cm}p{4.5cm}}
\toprule
\textbf{Component} & \textbf{Primary} & \textbf{Contributor} \\
\midrule
Data Ingestion (5 sources) & Aye Khin Khin Hpone & Dechathon Niamsa-Ard \\
Preprocessing \& Chunking & Dechathon Niamsa-Ard & Aye Khin Khin Hpone \\
Embedding \& Storage & Dechathon Niamsa-Ard & Aye Khin Khin Hpone \\
Hybrid Retrieval \& RAG & Aye Khin Khin Hpone & Dechathon Niamsa-Ard \\
LLM Backend Fallback & Dechathon Niamsa-Ard & Aye Khin Khin Hpone \\
Drift Detection \& Monitoring & Aye Khin Khin Hpone & Dechathon Niamsa-Ard \\
Infrastructure (SAM/CI-CD) & Dechathon Niamsa-Ard & Aye Khin Khin Hpone \\
Evaluation Framework & Aye Khin Khin Hpone & Dechathon Niamsa-Ard \\
Frontend (React) & Dechathon Niamsa-Ard & Aye Khin Khin Hpone \\
Unit Tests (49 tests) & Both & --- \\
Documentation \& Report & Both & --- \\
\bottomrule
\end{tabular}
\end{table}

Both team members participated in architectural design decisions, code review, and debugging throughout the project.

\section{Future Extensions}
\label{sec:future}

Several extensions are identified for post-prototype development:

\begin{itemize}[leftmargin=0.5in]
  \item Migration to Aurora Serverless if deployed outside Learner Lab constraints
  \item Expansion of data sources (StackExchange, academic conference proceedings)
  \item Cross-encoder reranking evaluation for precision-critical use cases
  \item Embedding model fine-tuning or replacement based on drift analysis findings
  \item Integration testing suite with containerized database and mocked AWS services
  \item User feedback loop for continuous retrieval quality improvement
  \item Source diversity balancing --- per-source ingestion quotas or source-aware reranking weights to prevent any single source from dominating retrieval results
  \item CloudFront HTTPS distribution --- serve the React frontend over HTTPS via CloudFront, eliminating browser ``Not secure'' warnings and improving production-readiness for demonstration
\end{itemize}

% ============================================================
% REFERENCES
% ============================================================
\chapter*{References}
\addcontentsline{toc}{chapter}{REFERENCES}

{\raggedright
\setlength{\parskip}{8pt}
\everypar{\hangindent=0.5in}

AWS. (2023). \textit{AWS well-architected framework}. Amazon Web Services. \url{https://aws.amazon.com/architecture/well-architected/}

Bayram, F., Ahmed, B. S., \& Kassler, A. (2022). From concept drift to model degradation: An overview on performance-aware drift detectors. \textit{Knowledge-Based Systems}, \textit{245}, 108632. \url{https://doi.org/10.1016/j.knosys.2022.108632}

Borgeaud, S., Mensch, A., Hoffmann, J., Cai, T., Rutherford, E., Millican, K., van den Driessche, G., Lespiau, J.-B., Damoc, B., Clark, A., de Las Casas, D., Guy, A., Menick, J., Ring, R., Hennigan, T., Huang, S., Maggiore, L., Jones, C., Cassirer, A., \ldots Sifre, L. (2022). Improving language models by retrieving from trillions of tokens. In \textit{Proceedings of the 39th International Conference on Machine Learning} (Vol. 162, pp. 2206--2240). PMLR.

Breck, E., Polyzotis, N., Roy, S., Whang, S., \& Zinkevich, M. (2017). The ML test score: A rubric for ML production readiness and technical debt reduction. In \textit{Proceedings of the IEEE International Conference on Big Data} (pp. 1123--1132). IEEE.

Brown, T., Mann, B., Ryder, N., Subbiah, M., Kaplan, J. D., Dhariwal, P., Neelakantan, A., Shyam, P., Sastry, G., Askell, A., Agarwal, S., Herbert-Voss, A., Krueger, G., Henighan, T., Child, R., Ramesh, A., Ziegler, D., Wu, J., Winter, C., \ldots Amodei, D. (2020). Language models are few-shot learners. \textit{Advances in Neural Information Processing Systems}, \textit{33}, 1877--1901.

Es, S., James, J., Anke, L. E., \& Schockaert, S. (2023). RAGAS: Automated evaluation of retrieval augmented generation. \textit{arXiv preprint arXiv:2309.15217}.

Gama, J., \v{Z}liobait\.{e}, I., Bifet, A., Pechenizkiy, M., \& Bouchachia, A. (2014). A survey on concept drift adaptation. \textit{ACM Computing Surveys}, \textit{46}(4), 1--37. \url{https://doi.org/10.1145/2523813}

Gao, Y., Xiong, Y., Gao, X., Jia, K., Pan, J., Bi, Y., Dai, Y., Sun, J., Wang, M., \& Wang, H. (2023). Retrieval-augmented generation for large language models: A survey. \textit{arXiv preprint arXiv:2312.10997}.

Izacard, G., Lewis, P., Lomeli, M., Hosseini, L., Petroni, F., Schick, T., Dwivedi-Yu, J., Joulin, A., Riedel, S., \& Grave, E. (2022). Atlas: Few-shot learning with retrieval augmented language models. \textit{Advances in Neural Information Processing Systems}, \textit{35}, 22251--22265.

Ji, Z., Lee, N., Frieske, R., Yu, T., Su, D., Xu, Y., Ishii, E., Bang, Y. J., Madotto, A., \& Fung, P. (2023). Survey of hallucination in natural language generation. \textit{ACM Computing Surveys}, \textit{55}(12), 1--38. \url{https://doi.org/10.1145/3571730}

Lewis, P., Perez, E., Piktus, A., Petroni, F., Karpukhin, V., Goyal, N., K{\"u}ttler, H., Lewis, M., Yih, W.-T., Rockt{\"a}schel, T., Riedel, S., \& Kiela, D. (2020). Retrieval-augmented generation for knowledge-intensive NLP tasks. \textit{Advances in Neural Information Processing Systems}, \textit{33}, 9459--9474.

Lu, J., Liu, A., Dong, F., Gu, F., Gama, J., \& Zhang, G. (2018). Learning under concept drift: A review. \textit{IEEE Transactions on Knowledge and Data Engineering}, \textit{31}(12), 2346--2363. \url{https://doi.org/10.1109/TKDE.2018.2876857}

Manakul, P., Liusie, A., \& Gales, M. J. F. (2023). SelfCheckGPT: Zero-resource black-box hallucination detection for generative large language models. In \textit{Proceedings of the 2023 Conference on Empirical Methods in Natural Language Processing} (pp. 9004--9017). Association for Computational Linguistics.

Min, S., Krishna, K., Lyu, X., Lewis, M., Yih, W.-T., Koh, P. W., Iyyer, M., Zettlemoyer, L., \& Hajishirzi, H. (2023). FActScoring: Fine-grained atomic evaluation of factual precision in long form text generation. In \textit{Proceedings of the 2023 Conference on Empirical Methods in Natural Language Processing} (pp. 12076--12100). Association for Computational Linguistics.

Nogueira, R., \& Cho, K. (2019). Passage re-ranking with BERT. \textit{arXiv preprint arXiv:1901.04085}.

Saad-Falcon, J., Khattab, O., Potts, C., \& Zaharia, M. (2023). ARES: An automated evaluation framework for retrieval-augmented generation systems. \textit{arXiv preprint arXiv:2311.09476}.

Sculley, D., Holt, G., Golovin, D., Davydov, E., Phillips, T., Ebner, D., Chaudhary, V., Young, M., Crespo, J.-F., \& Dennison, D. (2015). Hidden technical debt in machine learning systems. \textit{Advances in Neural Information Processing Systems}, \textit{28}, 2503--2511.

Shi, F., Chen, X., Misra, K., Scales, N., Dohan, D., Chi, E. H., Sch\"{a}rli, N., \& Zhou, D. (2023). Large language models can be easily distracted by irrelevant context. In \textit{Proceedings of the 40th International Conference on Machine Learning} (pp. 31210--31227). PMLR.

Thakur, N., Reimers, N., Daxenberger, J., \& Gurevych, I. (2021). BEIR: A heterogeneous benchmark for zero-shot evaluation of information retrieval models. In \textit{Proceedings of the 35th Conference on Neural Information Processing Systems---Datasets and Benchmarks Track}.

Wang, W., Wei, F., Dong, L., Bao, H., Yang, N., \& Zhou, M. (2020). MiniLM: Deep self-attention distillation for task-agnostic compression of pre-trained transformers. \textit{Advances in Neural Information Processing Systems}, \textit{33}, 5776--5788.

Zhang, T., Chen, Y., \& Zhou, Z. (2020). Machine learning operations (MLOps): Overview, definition, and architecture. \textit{IEEE Software}, \textit{37}(6), 1--17.

Zhang, Y., Chen, X., Zhang, Y., Hu, X., Liu, Z., \& Tang, J. (2024). A comprehensive survey on retrieval-augmented generation. \textit{ACM Transactions on Information Systems}. Advance online publication. \url{https://doi.org/10.1145/3703155}


}


% ============================================================
% APPENDIX A: SOURCE CODE AVAILABILITY
% ============================================================

\appendix

\chapter{Source Code Availability}

The implementation of \textit{TechPulse: Hybrid Retrieval-Augmented Generation System for Emerging Technology Intelligence} has been developed using Python 3.11 and modern cloud-native technologies.

To support reproducibility, transparency, and collaborative development, the project source code is maintained in a GitHub organization repository:

\begin{center}
\url{https://github.com/orgs/EZ2-555555/repositories}
\end{center}

The repository includes the following implemented components:

\begin{itemize}

\item Five data ingestion modules for collecting data from ArXiv API, Hacker News API, DEV.to API, GitHub Trending API, and curated RSS/Atom feeds

\item A 17-step preprocessing pipeline for comprehensive text normalization and token-based chunking using tiktoken (\texttt{cl100k\_base})

\item Semantic embedding pipeline using the \texttt{all-MiniLM-L6-v2} sentence-transformer model (384-dimensional embeddings)

\item Hybrid retrieval implementation including metadata filtering, cosine similarity search via pgvector, and reranking-lite scoring with configurable weights

\item RAG orchestration with three-layer hallucination verification (prompt grounding, RAGAS faithfulness, citation traceability)

\item Automatic LLM backend fallback chain supporting Groq (primary), Bedrock, Ollama, and HuggingFace Inference API

\item Dual-criteria drift detection combining Shewhart $3\sigma$ control charts with relative-drop thresholds

\item Infrastructure-as-Code definitions using AWS SAM for Learner Lab deployment (15+ AWS resources)

\item React-based frontend interface for query submission and citation-grounded response display

\item Evaluation framework with 25 curated queries, RAGAS metrics, and grid search optimization

\item Four-stage CI/CD pipeline using GitHub Actions (lint, test, build, deploy)

\item 49 unit tests across 5 test modules with full pass status

\item Docker Compose configuration for local development with 6 services

\end{itemize}

\chapter{Repository Structure}

The project repository is organized as follows:

\begin{lstlisting}
data-pipeline/
  src/
    config.py              # Centralized configuration (30+ settings)
    scheduler.py           # EventBridge schedule definitions
    api/main.py            # FastAPI + Mangum (/health, /ask, /drift)
    db/
      connection.py        # PostgreSQL connection pool
      init_schema.py       # Schema + pgvector + HNSW index
    embedding/embedder.py  # MiniLM singleton embedder
    ingestion/
      arxiv_ingester.py    # ArXiv API (cs.AI, cs.LG, cs.CL, cs.IR, cs.SE)
      hn_ingester.py       # Hacker News Firebase API
      devto_ingester.py    # DEV.to REST API (21 tags)
      github_ingester.py   # GitHub Search API (10 topics)
      rss_ingester.py      # 10 curated RSS/Atom feeds
    observability/
      __init__.py          # Deep health check (4 probes)
      drift.py             # Dual-criteria drift detection
    orchestrator/
      llm_backends.py      # 3-backend auto-fallback chain
      rag.py               # RAG pipeline + hallucination verification
    pipeline/run_pipeline.py  # End-to-end pipeline orchestrator
    preprocessing/chunker.py  # 17-step normalize + token chunking
    retrieval/retriever.py    # Baseline + hybrid 3-stage retrieval
    queue/                    # SQS producer/consumer
    storage/                  # S3 medallion operations
  tests/                # 49 unit tests (5 modules)
  evaluation/           # RAGAS evaluation + grid search
  infra/
    template-learnerlab.yaml  # SAM template (Learner Lab)
    samconfig-learnerlab.toml # SAM deployment config
  frontend/             # React SPA
  docker-compose.yml    # 6 local services
  .github/workflows/ci.yml   # 4-stage CI/CD
\end{lstlisting}

\chapter{System Demonstration Screenshots}
\label{app:screenshots}

This appendix provides supplementary screenshots demonstrating additional system functionality.

\section*{Health Check Endpoint}

Figure~\ref{fig:demo_health} shows the \texttt{/health} API endpoint response, confirming connectivity to all four service dependencies (PostgreSQL, S3, SQS, and LLM backend).

\begin{figure}[H]
\centering
\fbox{\parbox{0.9\textwidth}{\centering\vspace{2.5cm}
\textit{Screenshot: Browser or curl output of \texttt{GET /health} showing DB, S3, SQS, and LLM status.}
\\[4pt]
\textbf{File to place:} \texttt{figures/demo\_health.png}
\vspace{2.5cm}}}
\caption{Health check endpoint (\texttt{/health}): JSON response confirming database connectivity, document/chunk counts, S3 reachability, SQS reachability, and LLM backend status.}
\label{fig:demo_health}
\end{figure}

\section*{Drift Detection Endpoint}

Figure~\ref{fig:demo_drift} shows the \texttt{/drift} API endpoint returning a drift detection report, demonstrating the dual-criteria Shewhart $3\sigma$ monitoring mechanism.

\begin{figure}[H]
\centering
\fbox{\parbox{0.9\textwidth}{\centering\vspace{2.5cm}
\textit{Screenshot: Browser or curl output of \texttt{GET /drift} showing mean similarity, alert status, and probe results.}
\\[4pt]
\textbf{File to place:} \texttt{figures/demo\_drift.png}
\vspace{2.5cm}}}
\caption{Drift detection endpoint (\texttt{/drift}): JSON response showing probe query mean similarity, standard deviation, alert status, and Shewhart $3\sigma$ control analysis.}
\label{fig:demo_drift}
\end{figure}

\section*{API: RAG Query Response}

Figure~\ref{fig:demo_api} shows the raw JSON response from the \texttt{/ask} endpoint, displaying the structured output including the generated answer, source metadata, similarity scores, and retrieval mode.

\begin{figure}[H]
\centering
\fbox{\parbox{0.9\textwidth}{\centering\vspace{2.5cm}
\textit{Screenshot: Browser or curl output of \texttt{POST /ask} showing full JSON response with answer, sources, and scores.}
\\[4pt]
\textbf{File to place:} \texttt{figures/demo\_api.png}
\vspace{2.5cm}}}
\caption{RAG query endpoint (\texttt{POST /ask}): full JSON response showing the generated answer, retrieved source documents with titles, URLs, similarity scores, and retrieval mode metadata.}
\label{fig:demo_api}
\end{figure}

\end{document}
